%%%%%%%%%%%%%%%%%%%%%%%%%%%%%%%%%%%%%%%%%%%%%%%%%%%%%%%%%%%%%%%%%%%%%%%%%%%%%
%% Contact Maps in Systems Biology:
%% Individuation, Invariance, and the Derivation of Cellular State
%% from Partition Topology
%%
%% Author: Kundai Farai Sachikonye
%%         Technical University of Munich, School of Life Sciences
%%%%%%%%%%%%%%%%%%%%%%%%%%%%%%%%%%%%%%%%%%%%%%%%%%%%%%%%%%%%%%%%%%%%%%%%%%%%%

\documentclass[11pt,twocolumn]{article}

\usepackage[a4paper, top=2.5cm, bottom=2.5cm, left=2cm, right=2cm,
            columnsep=0.8cm]{geometry}
\usepackage{amsmath}
\usepackage{amsthm}
\usepackage{amssymb}
\usepackage{amsfonts}
\usepackage{mathtools}
\usepackage[ruled,vlined,linesnumbered]{algorithm2e}
\usepackage[colorlinks=true, linkcolor=blue, citecolor=blue, urlcolor=blue]{hyperref}
\usepackage[backend=bibtex, style=numeric, sorting=none]{biblatex}
\addbibresource{references.bib}
\usepackage{microtype}
\usepackage{graphicx}
\usepackage{booktabs}
\usepackage{enumitem}
\usepackage{xcolor}
\usepackage{thmtools}

%% Theorem environments
\theoremstyle{plain}
\newtheorem{theorem}{Theorem}[section]
\newtheorem{lemma}[theorem]{Lemma}
\newtheorem{proposition}[theorem]{Proposition}
\newtheorem{corollary}[theorem]{Corollary}

\theoremstyle{definition}
\newtheorem{definition}[theorem]{Definition}
\newtheorem{example}[theorem]{Example}
\newtheorem{remark}[theorem]{Remark}
\newtheorem{axiom}{Axiom}
\newtheorem{postulate}{Postulate}

%% Notation macros
\newcommand{\CC}{\mathcal{C}}
\newcommand{\GG}{\mathcal{G}}
\newcommand{\PP}{\mathcal{P}}
\newcommand{\NN}{\mathcal{N}}
\newcommand{\MM}{\mathcal{M}}
\newcommand{\FF}{\mathcal{F}}
\newcommand{\EE}{\mathcal{E}}
\newcommand{\HH}{\mathcal{H}}
\newcommand{\RR}{\mathcal{R}}
\newcommand{\Sk}{S_k}
\newcommand{\St}{S_t}
\newcommand{\Se}{S_e}
\newcommand{\bS}{\mathbf{S}}
\newcommand{\dS}{d_{\mathcal{S}}}
\newcommand{\betamin}{\beta_{\min}}
\newcommand{\reals}{\mathbb{R}}
\newcommand{\naturals}{\mathbb{N}}
\newcommand{\eps}{\varepsilon}
\newcommand{\norm}[1]{\left\lVert #1 \right\rVert}
\newcommand{\abs}[1]{\left\lvert #1 \right\rvert}

\title{On the Propagation of Partition Topology in Finite Graphs:\\ Minimum Sufficient Contact Maps in Systems Biology:}

\author{Kundai Farai Sachikonye\\[2pt]
\small Technical University of Munich \\
\small Chair of Experimental Bioinformatics\\
\small Maximus-von-Imhof-Forum 3, 85354 Freising, Germany\\
\small \texttt{kundai.sachikonye@wzw.tum.de}}

\date{\today}

\begin{document}

\maketitle

%%%%%%%%%%%%%%%%%%%%%%%%%%%%%%%%%%%%%%%%%%%%%%%%%%%%%%%%%%%%%%%%%%%%%%%%%%%%%
\begin{abstract}
%%%%%%%%%%%%%%%%%%%%%%%%%%%%%%%%%%%%%%%%%%%%%%%%%%%%%%%%%%%%%%%%%%%%%%%%%%%%%

We introduce a mathematical framework for representing cellular biochemical
networks as \emph{contact maps}: weighted graphs in which vertices are molecular
species and edges encode the minimum achievable thermodynamic separation between
species pairs. The framework rests on three postulates about how molecular
species are individuated within the crowded interior of a living cell: (1) a
species exists by being distinguishable from every other species present
(\emph{individuation by negation}); (2) the act of distinguishing two species
requires finite thermodynamic cost bounded below by an irreducible residue
$\beta > 0$; (3) the residue is monotone --- once deposited, it cannot be
annihilated without violating the second law.

From these postulates we derive a theory of \emph{contact} --- the state of
minimum inter-species partition depth --- and prove that contact is a
topological invariant of the biochemical network: it is preserved under all
physically realisable perturbations (enzyme inhibitions, concentration changes,
compartment rearrangements) and is irreducible in the sense that no coarser
topological fact carries equivalent information. We define a three-dimensional
\emph{S-entropy state space} $\mathcal{S} = [0,1]^3$ with coordinates
$(\Sk, \St, \Se)$ measuring kinetic flux entropy, topological connectivity
entropy, and energetic chemical potential entropy of each species. The
S-entropy distance between contacting species defines the contact map, and we
prove it is a well-defined metric on the space of distinguishable species.

We establish three main results: (i) the \emph{Contact Invariance Theorem},
showing that contact relationships are preserved under network perturbation and
partition refinement; (ii) the \emph{Contact Irreducibility Theorem}, proving
that contact is the minimum sufficient topological invariant of the network;
and (iii) the \emph{Contact Completeness Theorem}, establishing that the
contact map is a complete invariant of the network's partition topology --- two
networks with isomorphic contact maps have identical partition structure.

We demonstrate that the contact state is the unique privileged state of a
biochemical network from which all other states can be derived by partition
depth inflation (perturbation). Disease states, drug responses, and metabolic
shifts correspond to departures from contact, measurable as increased partition
depth along specific edges. We present algorithms for computing contact maps
from pathway databases, for backward navigation through the causal chain of
perturbations, and for $\ell^1$-optimal restoration of the contact state. The
framework is applied to glycolysis, the TCA cycle, oxidative phosphorylation,
and EGFR/MAPK signalling, with explicit computation of S-entropy coordinates,
contact costs, perturbation responses, and restoration targets. All quantities
are derived from publicly available thermodynamic and kinetic data (KEGG,
Reactome, HMDB, BRENDA) with zero free parameters.

\end{abstract}

%%%%%%%%%%%%%%%%%%%%%%%%%%%%%%%%%%%%%%%%%%%%%%%%%%%%%%%%%%%%%%%%%%%%%%%%%%%%%
\section{Introduction}
\label{sec:introduction}
%%%%%%%%%%%%%%%%%%%%%%%%%%%%%%%%%%%%%%%%%%%%%%%%%%%%%%%%%%%%%%%%%%%%%%%%%%%%%

The living cell presents a formidable challenge to mathematical representation.
Thousands of molecular species interact through tens of thousands of reactions
in a spatially compartmentalised, thermodynamically open, far-from-equilibrium
system \cite{kitano2002systems, palsson2015systems}. The dominant paradigms for
modelling such systems --- ordinary differential equations (ODEs) for reaction
kinetics \cite{klipp2016systems, tyson2003sniffers}, flux balance analysis (FBA)
for metabolic steady states \cite{orth2010flux, schellenberger2011quantitative},
and network topology analysis for structural properties \cite{barabasi2004network,
jeong2000large} --- have each illuminated essential aspects of cellular function.
Yet each also carries fundamental limitations.

ODE models require kinetic parameters (rate constants, Michaelis constants,
Hill coefficients) that are difficult to measure \emph{in vivo} and whose
values are context-dependent \cite{cornishbowden2012fundamentals, schomburg2004brenda}.
FBA models discard kinetics entirely, capturing only the stoichiometric skeleton
of the network at steady state \cite{orth2010flux}. Network topology analyses
abstract away both kinetics and thermodynamics, treating the network as a purely
combinatorial object \cite{barabasi2004network, ravasz2002hierarchical}.

In each case, the relationship between the mathematical representation and the
physical reality of the cell is unclear. An ODE model specifies a trajectory
through concentration space, but does not explain why certain trajectories are
accessible and others are not. An FBA model identifies feasible flux
distributions, but does not explain why certain distributions are preferred.
A network topology identifies hubs and motifs \cite{alon2007network}, but does
not explain why these structural features persist under perturbation.

\subsection{The Question of Invariance}

The present paper asks a different question: \emph{what is invariant?} When a
cell is perturbed --- by a drug, a mutation, a change in nutrient supply --- its
concentrations shift, its fluxes reroute, its entropy coordinates change. Yet
the cell remains \emph{the same cell}. What mathematical object captures the
identity of the cellular network that persists through all these changes?

We propose that the answer is the \emph{contact map}: the graph of minimum
thermodynamic separations between molecular species. Contact is the state of
maximum mutual constraint between two species --- the configuration in which
each species is most efficiently distinguished from the other. We prove that
contact relationships are topologically invariant: they cannot be broken by any
physically realisable perturbation, only deepened. The contact map is therefore
the skeleton of the network that persists through all accessible states.

\subsection{Individuation by Negation}

The philosophical starting point is the observation that a molecular species in
a cell exists not by virtue of its intrinsic properties alone, but by the
totality of what it is \emph{not}. Pyruvate in the cytoplasm is individuated
from glucose, from lactate, from alanine, from the thousands of other metabolites
in the cellular milieu. This individuation is not abstract: it requires physical
boundaries (different chemical potentials, different compartments, different
enzymatic contexts) that cost finite thermodynamic resources to establish and
maintain.

We formalise this observation as the principle of \emph{individuation by negation}:
a species $\mathcal{I}$ is defined by the systematic negation of everything it is
not, denoted $\bar{\mathcal{I}}$, and this negation has irreducible positive cost.
The cost is the \emph{partition depth} between $\mathcal{I}$ and $\bar{\mathcal{I}}$,
and its minimum value $\beta > 0$ is the \emph{irreducible residue} at the boundary
between the species and its complement.

\subsection{Contact as Ground State}

The minimum partition depth between two species --- the state where $\beta =
\betamin$ --- is \emph{contact}. We prove that contact is the ground state of
the pairwise interaction: all other configurations have higher partition depth,
and partition depth can only increase (or remain constant) under physical
operations. Contact is therefore the state from which all other states can be
derived by adding partition depth, but to which no operation can descend below.

This has a striking consequence for systems biology: the contact map of a
biochemical network is the unique privileged state from which all other states
(healthy, diseased, drug-treated, nutrient-shifted) can be derived. A disease is
not a different network; it is the \emph{same} contact map with elevated
partition depth along specific edges. The task of diagnosis is to identify which
edges have departed from contact, and the task of therapy is to find the
sparsest perturbation that restores them.

\subsection{Contributions}

Our main contributions are:

\begin{enumerate}[label=(\roman*)]
\item A rigorous axiomatic foundation for individuation by negation in
  biochemical networks, yielding the positivity of the irreducible residue
  and the non-instantaneity of species individuation
  (Section~\ref{sec:foundations}).

\item The theory of contact as a topological invariant: Contact Invariance
  (Theorem~\ref{thm:contact_invariance}), Contact Irreducibility
  (Theorem~\ref{thm:contact_irreducibility}), and Contact Completeness
  (Theorem~\ref{thm:contact_completeness}) (Section~\ref{sec:contact}).

\item The S-entropy state space formalism with kinetic, topological, and
  energetic entropy coordinates, and proof that the S-entropy distance is a
  well-defined metric on the space of distinguishable species
  (Section~\ref{sec:sentropy}).

\item The contact map construction algorithm from pathway database data, with
  explicit computation from public thermodynamic and kinetic databases
  (Section~\ref{sec:construction}).

\item Perturbation theory on contact maps: disease as partition depth inflation,
  backward navigation, and $\ell^1$-optimal restoration
  (Section~\ref{sec:perturbation}).

\item Application to four canonical pathways: glycolysis, TCA cycle, oxidative
  phosphorylation, and EGFR/MAPK signalling, with full numerical results
  (Section~\ref{sec:applications}).
\end{enumerate}

\subsection{Organisation}

Section~\ref{sec:foundations} establishes the axiomatic foundations.
Section~\ref{sec:contact} develops the theory of contact and proves the main
invariance theorems. Section~\ref{sec:sentropy} defines the S-entropy state
space. Section~\ref{sec:construction} gives the contact map construction
algorithm. Section~\ref{sec:perturbation} develops perturbation theory.
Section~\ref{sec:applications} presents applications. Section~\ref{sec:properties}
collects additional theoretical properties. Section~\ref{sec:discussion}
discusses implications. Section~\ref{sec:conclusion} concludes.


%%%%%%%%%%%%%%%%%%%%%%%%%%%%%%%%%%%%%%%%%%%%%%%%%%%%%%%%%%%%%%%%%%%%%%%%%%%%%
\section{Axiomatic Foundations}
\label{sec:foundations}
%%%%%%%%%%%%%%%%%%%%%%%%%%%%%%%%%%%%%%%%%%%%%%%%%%%%%%%%%%%%%%%%%%%%%%%%%%%%%

We begin with three postulates that formalise the physical constraints on the
individuation of molecular species in a living cell. These postulates are
deliberately minimal: each asserts a single physical fact that is
thermodynamically necessary, and together they generate the entire theory.

\subsection{The Three Postulates}

\begin{postulate}[Individuation by Negation]
\label{post:negation}
A molecular species $\mathcal{I}$ in a biochemical network exists if and only if
it is distinguishable from every other species in the network. Formally, let
$\Omega = \{\mathcal{I}_1, \ldots, \mathcal{I}_N\}$ be the set of molecular
species in a cell. Species $\mathcal{I}_i$ is \emph{individuated} if and only if
\[
  \mathcal{I}_i \neq \bar{\mathcal{I}}_i
  \quad\text{where}\quad
  \bar{\mathcal{I}}_i = \Omega \setminus \{\mathcal{I}_i\}
\]
and this distinction is established by at least one measurable physical
property (chemical potential, concentration, compartment, charge state, or
conformational state) that separates $\mathcal{I}_i$ from every member of
$\bar{\mathcal{I}}_i$.
\end{postulate}

\begin{remark}
Postulate~\ref{post:negation} is not the set-theoretic tautology $A \neq
\Omega \setminus A$. Its content is that species identity is \emph{constituted}
by the act of distinction: pyruvate is pyruvate because it is not glucose, not
lactate, not alanine, and not any of the other species present. Remove any one
of these distinctions (e.g., by making pyruvate and lactate indistinguishable)
and the identity of pyruvate is altered. This is the biological counterpart of
the philosophical principle \emph{omnis determinatio est negatio}.
\end{remark}

\begin{postulate}[Finite Individuation Cost]
\label{post:cost}
Every act of individuating species $\mathcal{I}_i$ from species $\mathcal{I}_j$
requires a thermodynamic cost bounded below by an irreducible residue
$\beta_{ij} > 0$. Formally, the \emph{partition depth} between
$\mathcal{I}_i$ and $\mathcal{I}_j$ satisfies
\[
  D(\mathcal{I}_i, \mathcal{I}_j) \geq \beta_{ij} > 0
\]
for all physically realisable states of the network. The residue $\beta_{ij}$
is the minimum thermodynamic cost of maintaining the distinction between
$\mathcal{I}_i$ and $\mathcal{I}_j$ and cannot be reduced to zero without
merging the two species into one.
\end{postulate}

\begin{remark}
The positivity of $\beta_{ij}$ has a concrete thermodynamic interpretation:
maintaining the distinction between two chemical species requires maintaining
a difference in at least one intensive variable (chemical potential, temperature,
pressure, or electrochemical potential), and any such difference in a
finite-volume system entails a positive free energy of separation
\cite{atkins2014physical, denbigh1981principles}. In the context of enzyme
kinetics, $\beta_{ij}$ corresponds to the minimum activation energy barrier
between two metabolites connected by an enzymatic reaction
\cite{cornishbowden2012fundamentals}.
\end{remark}

\begin{postulate}[Monotonicity of Partition Count]
\label{post:monotone}
The total partition count $M$ of the system --- the cumulative number of
individuation events that have occurred --- is strictly monotone increasing in
time:
\[
  t_2 > t_1 \implies M(t_2) > M(t_1).
\]
No physical operation can decrease $M$. Equivalently, the residues
$\{\beta_{ij}\}$ once deposited at partition boundaries cannot be annihilated;
they can only accumulate.
\end{postulate}

\begin{remark}
Postulate~\ref{post:monotone} is a restatement of the second law of
thermodynamics in the language of partitions: the entropy of the system
(which is the aggregate of unresolved residues whose specific causal links are
not tracked) can never decrease \cite{kondepudi2014modern, prigogine1967introduction}.
The partition count $M$ is a more fundamental quantity than entropy: entropy
arises when the partition count is coarse-grained over untracked causal links
\cite{shannon1948mathematical, jaynes1957information}.
\end{remark}

\subsection{Immediate Consequences}

\begin{definition}[Biochemical Partition Space]
\label{def:partition_space}
A \emph{biochemical partition space} is a tuple $(\Omega, D, \mu)$ where:
\begin{enumerate}[label=(\alph*)]
\item $\Omega = \{\mathcal{I}_1, \ldots, \mathcal{I}_N\}$ is a finite set of
  molecular species;
\item $D: \Omega \times \Omega \to \reals_{\geq 0}$ is a partition depth
  function satisfying $D(\mathcal{I}_i, \mathcal{I}_j) \geq \betamin > 0$ for
  all $i \neq j$ and $D(\mathcal{I}_i, \mathcal{I}_i) = 0$;
\item $\mu: \Omega \to \reals_{>0}$ assigns to each species a positive
  thermodynamic weight (its chemical potential, or more precisely, the absolute
  value of its Gibbs free energy of formation).
\end{enumerate}
\end{definition}

\begin{theorem}[Non-Instantaneity of Individuation]
\label{thm:noninstantaneous}
In any biochemical partition space $(\Omega, D, \mu)$, the individuation of any
species $\mathcal{I}_i$ from its complement $\bar{\mathcal{I}}_i$ requires at
least one committed partition step of cost $\geq \betamin > 0$.
\end{theorem}

\begin{proof}
Suppose, for contradiction, that species $\mathcal{I}_i$ can be individuated
from $\bar{\mathcal{I}}_i$ at zero cost. Then there exists at least one species
$\mathcal{I}_j \in \bar{\mathcal{I}}_i$ with $D(\mathcal{I}_i, \mathcal{I}_j)
= 0$, meaning $\mathcal{I}_i$ and $\mathcal{I}_j$ are indistinguishable. But
by Postulate~\ref{post:cost}, $D(\mathcal{I}_i, \mathcal{I}_j) \geq
\beta_{ij} > 0$, a contradiction. Therefore every individuation step costs at
least $\betamin = \min_{i \neq j} \beta_{ij} > 0$.

More constructively: since $\bar{\mathcal{I}}_i \neq \emptyset$ (there are at
least two species in any non-trivial network), there exists at least one
$\mathcal{I}_j \in \bar{\mathcal{I}}_i$ with $D(\mathcal{I}_i,
\mathcal{I}_j) \geq \betamin$. The separation of $\mathcal{I}_i$ from
$\mathcal{I}_j$ is a partition step of cost $\geq \betamin > 0$, establishing
the claimed lower bound. \qed
\end{proof}

\begin{definition}[Individuation Residue]
\label{def:residue}
Let $\mathcal{I}_i$ and $\mathcal{I}_j$ be two adjacent species in the network
(connected by at least one reaction). The \emph{individuation residue} of the
pair $(\mathcal{I}_i, \mathcal{I}_j)$ is the minimum achievable partition depth:
\[
  \beta(\mathcal{I}_i, \mathcal{I}_j)
  = \inf_{\text{all accessible states}} D(\mathcal{I}_i, \mathcal{I}_j).
\]
By Postulate~\ref{post:cost}, $\beta(\mathcal{I}_i, \mathcal{I}_j) > 0$.
By Postulate~\ref{post:monotone}, once the residue is deposited at the boundary,
it cannot be annihilated.
\end{definition}

\begin{proposition}[Residue Chain]
\label{prop:residue_chain}
The individuation of species $\mathcal{I}_i$ from $\mathcal{I}_j$ deposits a
residue $\beta_{ij}$ at the boundary between them. This residue is itself a
physical difference that modifies the boundary conditions of all neighbouring
species. Specifically, for any species $\mathcal{I}_k$ adjacent to
$\mathcal{I}_j$ but not previously adjacent to $\mathcal{I}_i$, the
individuation of $(\mathcal{I}_i, \mathcal{I}_j)$ may create a new adjacency
between $\mathcal{I}_i$ and $\mathcal{I}_k$, requiring a new individuation
step of cost $\geq \betamin$.
\end{proposition}

\begin{proof}
The partition depth $D(\mathcal{I}_i, \mathcal{I}_j) \geq \beta_{ij} > 0$
constitutes a measurable physical difference at the boundary between
$\mathcal{I}_i$ and $\mathcal{I}_j$ \cite{atkins2014physical}. Any species
$\mathcal{I}_k$ that shares a reaction with $\mathcal{I}_j$ is sensitive to the
boundary conditions of $\mathcal{I}_j$. If the residue $\beta_{ij}$ alters the
effective concentration, flux, or chemical potential profile of $\mathcal{I}_j$
as seen from $\mathcal{I}_k$, then the partition depth $D(\mathcal{I}_j,
\mathcal{I}_k)$ changes, and a new individuation step between $\mathcal{I}_i$
and $\mathcal{I}_k$ (mediated through $\mathcal{I}_j$) may become necessary.

Since the network is finite and connected, the propagation of residues through
the contact chain terminates after at most $N-1$ steps (spanning tree of the
network graph), where $N = |\Omega|$. Each step deposits a residue of cost
$\geq \betamin$, and by Postulate~\ref{post:monotone}, these residues accumulate
monotonically. \qed
\end{proof}

\begin{corollary}[Contact Chain Self-Reinforcement]
\label{cor:self_reinforce}
Adding a new individuation event to the network increases the definitional
precision of every species in the connected component. Specifically, if species
$\mathcal{I}_i$ and $\mathcal{I}_j$ are brought into contact, the partition
depth of every other pair $(\mathcal{I}_k, \mathcal{I}_l)$ in the same
connected component either increases or remains constant. No partition depth
decreases.
\end{corollary}

\begin{proof}
By Postulate~\ref{post:monotone}, the total partition count $M$ increases with
each new individuation event. The partition depth $D(\mathcal{I}_k,
\mathcal{I}_l)$ is a component of $M$ (it counts the individuation events
separating $\mathcal{I}_k$ from $\mathcal{I}_l$). Since $M$ is monotone
increasing and $D$ is a sum of non-negative contributions from each
individuation event, each $D(\mathcal{I}_k, \mathcal{I}_l)$ is monotone
non-decreasing. \qed
\end{proof}


%%%%%%%%%%%%%%%%%%%%%%%%%%%%%%%%%%%%%%%%%%%%%%%%%%%%%%%%%%%%%%%%%%%%%%%%%%%%%
\section{Contact as Topological Invariant}
\label{sec:contact}
%%%%%%%%%%%%%%%%%%%%%%%%%%%%%%%%%%%%%%%%%%%%%%%%%%%%%%%%%%%%%%%%%%%%%%%%%%%%%

We now formalise the central concept of \emph{contact} and prove its topological
invariance and irreducibility.

\subsection{Definitions}

\begin{definition}[Contact]
\label{def:contact}
Two molecular species $\mathcal{I}_i$ and $\mathcal{I}_j$ in a biochemical
partition space $(\Omega, D, \mu)$ are in \emph{contact} if their partition
depth achieves the minimum:
\[
  D(\mathcal{I}_i, \mathcal{I}_j)
  = \beta(\mathcal{I}_i, \mathcal{I}_j)
  = \betamin(\mathcal{I}_i, \mathcal{I}_j).
\]
We write $C(\mathcal{I}_i, \mathcal{I}_j) = 1$ when the pair is in contact,
and $C(\mathcal{I}_i, \mathcal{I}_j) = 0$ otherwise.
\end{definition}

\begin{remark}
Contact is not the same as adjacency. Two species can be adjacent in the
reaction network (connected by an enzyme-catalysed reaction) without being in
contact: their partition depth may exceed the minimum due to perturbations
(enzyme inhibition, substrate depletion, compartment disruption). Contact is
the \emph{ground state} of the adjacency: the configuration of minimum
thermodynamic separation consistent with the species remaining distinguishable.
\end{remark}

\begin{definition}[Contact Graph]
\label{def:contact_graph}
The \emph{contact graph} of a biochemical partition space $(\Omega, D, \mu)$
is the undirected weighted graph $G_C = (V, E, w)$ where:
\begin{itemize}
\item $V = \Omega$ (one vertex per species);
\item $E = \{(\mathcal{I}_i, \mathcal{I}_j) : \mathcal{I}_i \text{ and }
  \mathcal{I}_j \text{ share a reaction}\}$ (edges from the reaction network);
\item $w(\mathcal{I}_i, \mathcal{I}_j) = \beta(\mathcal{I}_i, \mathcal{I}_j)$
  (the irreducible residue is the edge weight).
\end{itemize}
\end{definition}

\begin{definition}[Contact Map]
\label{def:contact_map}
The \emph{contact map} of the biochemical partition space is the function
\[
  \CC: E \to \reals_{>0}, \qquad
  \CC(\mathcal{I}_i, \mathcal{I}_j) = \beta(\mathcal{I}_i, \mathcal{I}_j)
\]
assigning to each edge the minimum partition depth between its endpoints.
\end{definition}

\begin{definition}[Partition Depth Inflation]
\label{def:inflation}
A \emph{perturbation} of the biochemical partition space is a function
$\delta: E \to \reals_{\geq 0}$ that modifies the partition depth of each edge:
\[
  D'(\mathcal{I}_i, \mathcal{I}_j)
  = D(\mathcal{I}_i, \mathcal{I}_j) + \delta(\mathcal{I}_i, \mathcal{I}_j).
\]
Since $\delta \geq 0$ (by Postulate~\ref{post:monotone}, partition depth cannot
decrease), any perturbation is a \emph{partition depth inflation}: it can only
increase or maintain the partition depth along each edge.
\end{definition}

\begin{definition}[Contact Topology]
\label{def:contact_topology}
The \emph{contact topology} $\mathcal{T}_C$ on $V$ is the topology generated by
the contact graph: the open sets are unions of cliques in $G_C$ (sets of
mutually contacting species). Two species are in the same connected component
of $\mathcal{T}_C$ if and only if there exists a path of contact relationships
connecting them.
\end{definition}

\subsection{Main Theorems}

\begin{theorem}[Contact Invariance]
\label{thm:contact_invariance}
Let $(\Omega, D, \mu)$ be a biochemical partition space and let $\delta: E \to
\reals_{\geq 0}$ be any perturbation. If $C(\mathcal{I}_i, \mathcal{I}_j) = 1$
in the unperturbed state, then $C(\mathcal{I}_i, \mathcal{I}_j) = 1$ in the
perturbed state. That is, contact is preserved under all physically realisable
perturbations.
\end{theorem}

\begin{proof}
The proof has three parts: (a) contact is preserved under partition depth
inflation; (b) partition depth inflation is the only physically realisable
class of perturbation; (c) therefore contact is invariant.

\textbf{(a)} Suppose $C(\mathcal{I}_i, \mathcal{I}_j) = 1$, meaning
$D(\mathcal{I}_i, \mathcal{I}_j) = \beta(\mathcal{I}_i, \mathcal{I}_j)$. Under
perturbation $\delta$, the new partition depth is $D'(\mathcal{I}_i,
\mathcal{I}_j) = D(\mathcal{I}_i, \mathcal{I}_j) + \delta(\mathcal{I}_i,
\mathcal{I}_j)$. The new minimum achievable partition depth is $\beta'
(\mathcal{I}_i, \mathcal{I}_j) = \beta(\mathcal{I}_i, \mathcal{I}_j)$
(the irreducible residue is intrinsic to the species pair and does not change
under perturbation; it depends only on the chemical identities of the two
species, not on the state of the network).

We must show that the contact relationship $C$ is defined at the level of the
minimum, not the current value. The perturbation increases $D'$ above $\beta$,
so the pair is no longer at the minimum partition depth. However, the
\emph{contact relationship} asserts that the minimum is achievable: the pair
\emph{can} return to $D = \beta$ if the perturbation is removed. Since
$\beta > 0$ is an intrinsic property of the species pair (determined by their
chemical structures and the thermodynamics of their boundary), it is invariant
under perturbation.

More precisely: the contact graph $G_C$ records which pairs of species have
$\beta > 0$ (which is all adjacent pairs, by Postulate~\ref{post:cost}). A
perturbation changes the \emph{current} partition depth $D$ but does not
alter the \emph{minimum achievable} partition depth $\beta$. The edge set of
$G_C$ is therefore invariant under perturbation.

\textbf{(b)} By Postulate~\ref{post:monotone}, the only physically realisable
perturbations are those with $\delta \geq 0$ (partition depth can only increase
or stay constant). An operation that decreases $D$ below $\beta$ would require
annihilating the residue, which violates Postulate~\ref{post:monotone}. An
operation that removes an edge from $G_C$ would require making two species
indistinguishable ($D = 0$), which violates Postulate~\ref{post:cost}. Neither
is physically realisable.

\textbf{(c)} Combining (a) and (b): the edge set of the contact graph $G_C$
is invariant under all physically realisable perturbations. The edge weights
$w = \beta$ are intrinsic to the species pairs and are also invariant. Therefore
the contact map $\CC$ is invariant under all physically realisable
perturbations. \qed
\end{proof}

\begin{theorem}[Contact Irreducibility]
\label{thm:contact_irreducibility}
Contact is the minimum sufficient topological invariant of the biochemical
network. No strictly coarser topological fact about the network carries the
same information across all perturbation levels.
\end{theorem}

\begin{proof}
Let $F$ be any topological predicate on pairs of species that is:
\begin{enumerate}[label=(\roman*)]
\item \emph{invariant}: preserved under all physically realisable perturbations;
\item \emph{non-trivial}: there exist non-adjacent species $\mathcal{I}_i,
  \mathcal{I}_j$ with $F(\mathcal{I}_i, \mathcal{I}_j) = 0$.
\end{enumerate}

We claim that if $F$ is coarser than contact (i.e., $C(\mathcal{I}_i,
\mathcal{I}_j) = 1 \implies F(\mathcal{I}_i, \mathcal{I}_j) = 1$) and
$F$ is invariant, then $F$ is equivalent to contact.

Suppose $F(\mathcal{I}_i, \mathcal{I}_j) = 1$ but $C(\mathcal{I}_i,
\mathcal{I}_j) = 0$ for some pair, meaning $\mathcal{I}_i$ and $\mathcal{I}_j$
are not adjacent in the reaction network. Since they are not adjacent, there is
no reaction connecting them, and hence no partition boundary between them. The
partition depth $D(\mathcal{I}_i, \mathcal{I}_j)$ is not defined directly but
only through paths in the contact graph. Any topological predicate that assigns
a positive value to a non-adjacent pair must be derivable from the path structure
of $G_C$, which is itself determined by the edge set $E$ and the edge weights
$\beta$. But $E$ and $\beta$ are exactly the data of the contact map $\CC$.
Therefore $F$ is determined by $\CC$, and any predicate that agrees with
$\CC$ on all adjacent pairs and is determined by $\CC$ on non-adjacent pairs
is equivalent to $\CC$ on the domain of non-trivial predicates.

If $F$ is strictly coarser than contact (i.e., $F(\mathcal{I}_i,
\mathcal{I}_j) = 1$ for some pair with $C(\mathcal{I}_i, \mathcal{I}_j) = 0$
and $F(\mathcal{I}_k, \mathcal{I}_l) = 0$ for some pair with
$C(\mathcal{I}_k, \mathcal{I}_l) = 0$), then $F$ does not distinguish certain
non-adjacent pairs that the contact map does distinguish via path structure.
Such an $F$ loses information about the network topology and is therefore not
sufficient for recovering the full partition structure.

We conclude that the only monotone, non-trivial, invariant topological predicate
that carries the same information as the contact map is the contact map itself.
Contact is irreducible. \qed
\end{proof}

\begin{theorem}[Contact Completeness]
\label{thm:contact_completeness}
The contact map $\CC: E \to \reals_{>0}$ is a complete invariant of the
contact topology: two biochemical partition spaces $(\Omega_1, D_1, \mu_1)$ and
$(\Omega_2, D_2, \mu_2)$ with isomorphic contact maps (as weighted graphs)
have isomorphic contact topologies.
\end{theorem}

\begin{proof}
The contact topology $\mathcal{T}_C$ is determined by the contact graph
$G_C = (V, E, w)$: the open sets of $\mathcal{T}_C$ are determined by $E$, and
the metric structure is determined by the weights $w = \beta$. If two networks
have isomorphic contact maps (meaning there exists a bijection $\phi: V_1 \to
V_2$ preserving both the edge set and the edge weights), then their contact
graphs are isomorphic as weighted graphs, and hence their contact topologies are
homeomorphic.

Conversely, if the contact topologies are homeomorphic, then the underlying
graphs are isomorphic (the topology determines the graph) and the weights
are preserved (the weights are topological invariants by
Theorem~\ref{thm:contact_invariance}).

Therefore the contact map is a complete invariant: it determines and is
determined by the contact topology. \qed
\end{proof}
\begin{figure*}[!htbp]
    \centering
    \includegraphics[width=\textwidth]{figures/panel_7_residue_compartments.png}
    \caption{\textbf{Residue chain propagation, compartmental structure, and triple coherence under perturbation.}
    (\textbf{A})~Total partition depth $\sum_e \beta(e)$ versus perturbation strength $\alpha$ (log scale) for glycolysis (blue) and TCA (orange). Downward triangles mark the minimum-depth configuration. Glycolysis shows a broad minimum near $\alpha \approx 1.0$--$1.5$ with total depth $\approx 4.13$, confirming the self-reinforcement principle: the contact state approximately minimises total partition depth. TCA shows a monotonically decreasing profile over the tested range, consistent with its cyclic topology redistributing perturbation energy.
    (\textbf{B})~Three-dimensional compartmental structure of EGFR/MAPK signalling. Vertices are positioned by $(\Sk, \St)$ and stratified along the $z$-axis by compartment (Nucleus = 0, Cytoplasm = 1, Membrane = 2, Extracellular = 3). Red edges cross compartment boundaries; grey edges remain within compartments. The four-compartment architecture of the MAPK cascade is recovered directly from the S-entropy embedding.
    (\textbf{C})~Compartmental factorisation (Proposition~7.2): mean intra-compartmental contact cost (green) versus mean inter-compartmental contact cost (red) for OxPhos and EGFR/MAPK. EGFR/MAPK satisfies the factorisation inequality $\bar{\mathcal{C}}_{\text{inter}} / \bar{\mathcal{C}}_{\text{intra}} = 8.9$, confirming that membrane-crossing edges carry substantially higher partition depth. OxPhos shows a smaller ratio with $\bar{\mathcal{C}}_{\text{inter}} < \bar{\mathcal{C}}_{\text{intra}}$, reflecting the tight energetic coupling across the inner mitochondrial membrane in oxidative phosphorylation.
    (\textbf{D})~Triple coherence $R(\alpha)$ response curves for all four pathways. Glycolysis (blue, $R = 0.463$ at $\alpha = 1$) maintains stable coherence across the perturbation range. TCA (orange) increases from $R = 0.111$ to $R = 0.172$ under perturbation, indicating that IDH2 inhibition paradoxically improves rank alignment. OxPhos (green, $R \approx -0.08$) shows weak anti-correlation throughout. EGFR/MAPK (red, $R = 0.152$) remains invariant, consistent with the near-unity visibility under KRAS G12V.}
    \label{fig:panel-residue-compartments}
\end{figure*}


%%%%%%%%%%%%%%%%%%%%%%%%%%%%%%%%%%%%%%%%%%%%%%%%%%%%%%%%%%%%%%%%%%%%%%%%%%%%%
\section{The S-Entropy State Space}
\label{sec:sentropy}
%%%%%%%%%%%%%%%%%%%%%%%%%%%%%%%%%%%%%%%%%%%%%%%%%%%%%%%%%%%%%%%%%%%%%%%%%%%%%

To quantify the partition depth between species and to define the contact map
computationally, we introduce a three-dimensional state space that captures the
thermodynamic character of each molecular species.

\subsection{Definitions}

\begin{definition}[S-Entropy Coordinates]
\label{def:sentropy_coords}
Let $\mathcal{I}_i$ be a molecular species in a biochemical partition space
$(\Omega, D, \mu)$ embedded in a reaction network $G = (V, E)$. The
\emph{S-entropy coordinates} of $\mathcal{I}_i$ are the triple
$\bS(\mathcal{I}_i) = (\Sk(\mathcal{I}_i), \St(\mathcal{I}_i),
\Se(\mathcal{I}_i)) \in [0,1]^3$ defined as follows.

\paragraph{Kinetic entropy $\Sk(\mathcal{I}_i)$} (normalised flux magnitude):
\begin{equation}\label{eq:Sk}
  \Sk(\mathcal{I}_i)
  = \frac{\displaystyle\sum_{j \in \mathcal{N}(i)} |J_{ij}|}
         {\displaystyle\max_{k \in V} \sum_{j \in \mathcal{N}(k)} |J_{kj}|}
  \in [0,1],
\end{equation}
where $J_{ij}$ is the net flux from $\mathcal{I}_i$ to $\mathcal{I}_j$
(in mol/s or equivalent units), $\mathcal{N}(i)$ is the set of neighbours of
$\mathcal{I}_i$ in the reaction network, and the denominator normalises by the
maximum total flux magnitude over all species.

\paragraph{Topological entropy $\St(\mathcal{I}_i)$} (normalised weighted degree):
\begin{equation}\label{eq:St}
  \St(\mathcal{I}_i)
  = \frac{\displaystyle\sum_{j \in \mathcal{N}(i)} G_{ij}}
         {\displaystyle\max_{k \in V} \sum_{j \in \mathcal{N}(k)} G_{kj}}
  \in [0,1],
\end{equation}
where $G_{ij}$ is the \emph{conductance} of the edge $(\mathcal{I}_i,
\mathcal{I}_j)$, defined as the proportionality constant between the flux
and the chemical potential difference:
\begin{equation}\label{eq:conductance}
  J_{ij} = G_{ij} \cdot (\mu_i - \mu_j),
\end{equation}
with $\mu_i$ the chemical potential of species $\mathcal{I}_i$
\cite{onsager1931reciprocal, kondepudi2014modern}.

\paragraph{Energetic entropy $\Se(\mathcal{I}_i)$} (normalised chemical potential):
\begin{equation}\label{eq:Se}
  \Se(\mathcal{I}_i)
  = \frac{\mu_i - \mu_{\min}}{\mu_{\max} - \mu_{\min}}
  \in [0,1],
\end{equation}
where $\mu_i = \mu_i^\circ + RT \ln[\mathcal{I}_i]$ is the chemical potential
of species $\mathcal{I}_i$ \cite{atkins2014physical, alberty2003thermodynamics},
$\mu_{\min} = \min_k \mu_k$, and $\mu_{\max} = \max_k \mu_k$.
\end{definition}

\begin{remark}[Independence of the Three Coordinates]
\label{rem:independence}
The three S-entropy coordinates capture genuinely independent aspects of the
species' role in the network:
\begin{itemize}
\item $\Sk$ measures \emph{kinetic activity}: how much flux passes through the
  species. A species with $\Sk = 1$ is the most kinetically active node in the
  network; $\Sk = 0$ is a kinetic dead end.
\item $\St$ measures \emph{topological connectivity}: how strongly coupled the
  species is to its neighbours via conductance. A species with $\St = 1$ is the
  most connected node; $\St = 0$ is topologically isolated.
\item $\Se$ measures \emph{energetic position}: where the species sits on the
  thermodynamic landscape. A species with $\Se = 1$ has the highest chemical
  potential (most energy-rich); $\Se = 0$ has the lowest (most stable).
\end{itemize}
There is no sum constraint $\Sk + \St + \Se = 1$; the three coordinates are
independent axes of a cube $[0,1]^3$. A glucose molecule at the start of
glycolysis might have $\bS = (0.7, 0.3, 0.95)$ --- high flux, moderate
connectivity, high chemical potential --- while pyruvate at the end might have
$\bS = (0.8, 0.5, 0.05)$ --- high flux, moderate connectivity, low chemical
potential.
\end{remark}

\begin{definition}[S-Entropy State Space]
\label{def:state_space}
The \emph{S-entropy state space} is $\mathcal{S} = [0,1]^3$ equipped with the
Euclidean metric:
\begin{equation}\label{eq:dS}
  \dS(\bS_1, \bS_2)
  = \sqrt{(\Sk^{(1)} - \Sk^{(2)})^2
        + (\St^{(1)} - \St^{(2)})^2
        + (\Se^{(1)} - \Se^{(2)})^2}.
\end{equation}
The maximum distance in $\mathcal{S}$ is $\sqrt{3}$ (between the vertices
$(0,0,0)$ and $(1,1,1)$ of the unit cube).
\end{definition}

\begin{definition}[S-Entropy Contact Map]
\label{def:sentropy_contact_map}
The \emph{S-entropy contact map} is the function
\begin{equation}\label{eq:CM}
  \CC(\mathcal{I}_i, \mathcal{I}_j)
  = \dS\bigl(\bS(\mathcal{I}_i), \bS(\mathcal{I}_j)\bigr)
\end{equation}
for every edge $(\mathcal{I}_i, \mathcal{I}_j) \in E$ of the reaction network.
\end{definition}

\subsection{Well-Definedness and Metric Properties}

\begin{lemma}[Well-Definedness of S-Entropy Coordinates]
\label{lem:well_defined}
Let $(\Omega, D, \mu)$ be a biochemical partition space with reaction network
$G = (V, E)$, and let each species $\mathcal{I}_i$ have well-defined chemical
potential $\mu_i$, concentration $[\mathcal{I}_i] > 0$, and at least one
reaction partner. Then the S-entropy coordinates $\bS(\mathcal{I}_i) =
(\Sk, \St, \Se) \in [0,1]^3$ are well-defined and finite.
\end{lemma}

\begin{proof}
\textbf{Kinetic entropy $\Sk$}: The flux $J_{ij} = G_{ij}(\mu_i - \mu_j)$ is
well-defined because $G_{ij} \geq 0$ (conductance is non-negative) and $\mu_i,
\mu_j$ are finite (by the finiteness of chemical potentials for species at
positive concentration \cite{atkins2014physical}). The sum $\sum_j |J_{ij}|$ is
finite because $|\mathcal{N}(i)|$ is finite. The denominator $\max_k \sum_j
|J_{kj}| > 0$ because the network has at least one edge with $G_{ij} > 0$ and
$\mu_i \neq \mu_j$ (otherwise the network would be at thermodynamic equilibrium
with zero flux everywhere, contradicting the far-from-equilibrium condition of
living cells). Therefore $\Sk \in [0,1]$ is well-defined.

\textbf{Topological entropy $\St$}: The conductance $G_{ij} \geq 0$ is
well-defined for each edge. The sum $\sum_j G_{ij}$ is finite. The denominator
$\max_k \sum_j G_{kj} > 0$ because at least one species has positive
conductance to at least one neighbour. Therefore $\St \in [0,1]$.

\textbf{Energetic entropy $\Se$}: The chemical potential $\mu_i = \mu_i^\circ
+ RT \ln[\mathcal{I}_i]$ is well-defined for $[\mathcal{I}_i] > 0$
\cite{atkins2014physical}. The denominator $\mu_{\max} - \mu_{\min} > 0$
because in a non-trivial network with at least two species at different chemical
potentials (otherwise no thermodynamic driving force exists). Therefore
$\Se \in [0,1]$. \qed
\end{proof}

\begin{figure*}[!htbp]
    \centering
    \includegraphics[width=\textwidth]{figures/panel_1_sentropy_embeddings.png}
    \caption{\textbf{S-entropy vertex embeddings $\mathbf{S}: V \to [0,1]^{3}$ for four canonical biochemical networks.}
    Each subplot displays the three-dimensional embedding of every vertex (molecular species) in the S-entropy state space $(\Sk, \St, \Se)$, with colour mapped to the energetic coordinate $\Se$ (viridis scale). Vertical drop lines project each vertex onto the $(\Sk, \St)$ floor to aid depth perception.
    (\textbf{A})~Glycolysis ($P_{10}$, 10 vertices): the preparatory-phase vertices (Glucose, G6P) occupy the high-$\Sk$ region while payoff-phase vertices (3PG, 2PG, Pyruvate) cluster near the origin, reflecting the asymmetric distribution of kinetic flux.
    (\textbf{B})~TCA cycle ($C_9$, 9 vertices): the cyclic topology distributes vertices more uniformly across the state space; Malate and Acetyl-CoA anchor the high-$\Sk$ extremes while Isocitrate occupies a low-flux position.
    (\textbf{C})~Oxidative phosphorylation (8 vertices, mixed graph): vertices stratify primarily along $\Se$, with electron carriers (NADH, CoQ, CytC) at high energetic entropy and the ATP synthesis node at $\Se = 0$.
    (\textbf{D})~EGFR/MAPK signalling (10 vertices, directed graph with feedback): the ligand--receptor pair (EGF, EGFR) dominates the high-$\Sk$, high-$\St$ corner while downstream effectors (MYC, CycD) collapse to near-zero kinetic and topological coordinates, consistent with signal amplification along the cascade.}
    \label{fig:panel-sentropy-embeddings}
\end{figure*}

\begin{theorem}[S-Entropy Distance is a Metric]
\label{thm:metric}
The function $\dS: \mathcal{S} \times \mathcal{S} \to \reals_{\geq 0}$ defined
by Equation~\eqref{eq:dS} is a metric on the space of distinguishable species.
\end{theorem}

\begin{proof}
Since $\dS$ is the restriction of the Euclidean metric on $\reals^3$ to the
compact subset $[0,1]^3$, it inherits all metric properties:
\begin{enumerate}[label=(\alph*)]
\item \emph{Non-negativity}: $\dS(\bS_1, \bS_2) \geq 0$ with equality iff
  $\bS_1 = \bS_2$.
\item \emph{Symmetry}: $\dS(\bS_1, \bS_2) = \dS(\bS_2, \bS_1)$.
\item \emph{Triangle inequality}: $\dS(\bS_1, \bS_3) \leq \dS(\bS_1, \bS_2)
  + \dS(\bS_2, \bS_3)$.
\end{enumerate}
For distinguishable species, $\bS(\mathcal{I}_i) \neq \bS(\mathcal{I}_j)$
whenever $i \neq j$ (otherwise the two species would be indistinguishable in
S-entropy space, contradicting Postulate~\ref{post:negation} which requires
each species to be distinguishable from all others). Therefore
$\dS(\bS(\mathcal{I}_i), \bS(\mathcal{I}_j)) > 0$ for $i \neq j$. \qed
\end{proof}

\begin{definition}[Partition Ground State]
\label{def:ground_state}
The \emph{partition ground state} of a species $\mathcal{I}_i$ is the point
$(1, 0, 0, s)$ in the partition state space $(n, \ell, m, s)$, corresponding to
$\Sk = 0$, $\St = 0$, $\Se = 1/2$ in S-entropy coordinates. The S-entropy
distance from the ground state measures how far a species is from the simplest
possible partition configuration.
\end{definition}

\begin{theorem}[S-Entropy as Partition Distance]
\label{thm:partition_distance}
The S-entropy vector $\bS(\mathcal{I}_i) = (\Sk, \St, \Se)$ measures the
distance of species $\mathcal{I}_i$ from the partition ground state
$(0, 0, 1/2)$. This distance is a well-defined metric on partition states.
\end{theorem}

\begin{proof}
The ground state has $\Sk = 0$ (no kinetic flux), $\St = 0$ (no topological
connectivity), and $\Se = 1/2$ (midpoint chemical potential). The Euclidean
distance from the ground state is:
\[
  d = \sqrt{\Sk^2 + \St^2 + (\Se - 1/2)^2}.
\]
This is non-negative, zero only at the ground state, symmetric (trivially), and
satisfies the triangle inequality (as a Euclidean distance). \qed
\end{proof}

\subsection{Triple Coherence}

\begin{definition}[Triple Coherence]
\label{def:triple_coherence}
The \emph{triple coherence} $R$ of a contact map is the Spearman rank
correlation \cite{spearman1904proof} among the three S-entropy coordinate
vectors $(\Sk^{(1)}, \ldots, \Sk^{(N)})$, $(\St^{(1)}, \ldots, \St^{(N)})$,
$(\Se^{(1)}, \ldots, \Se^{(N)})$:
\begin{equation}\label{eq:coherence}
  R = \frac{1}{3}\bigl(
    \rho(\Sk, \St) + \rho(\St, \Se) + \rho(\Sk, \Se)
  \bigr),
\end{equation}
where $\rho$ is the Spearman rank correlation coefficient. $R \in [-1, 1]$,
with $R = 1$ indicating perfect coherence (all three entropy measures agree on
the ordering of species) and $R = 0$ indicating no coherence.
\end{definition}

\begin{proposition}[Coherence at Contact]
\label{prop:coherence_contact}
At the contact state (all edges at minimum partition depth), the triple coherence
$R$ achieves its maximum value for the given network topology. Any perturbation
that increases partition depth along a subset of edges reduces $R$.
\end{proposition}

\begin{proof}
At contact, the flux pattern $\{J_{ij}\}$ is determined by the conductance
matrix $\{G_{ij}\}$ and the chemical potential vector $\{\mu_i\}$ via
$J_{ij} = G_{ij}(\mu_i - \mu_j)$. The three S-entropy coordinates are then
deterministic functions of $\{G_{ij}\}$ and $\{\mu_i\}$:
\begin{align}
  \Sk(\mathcal{I}_i) &\propto \sum_j G_{ij} |\mu_i - \mu_j|, \\
  \St(\mathcal{I}_i) &\propto \sum_j G_{ij}, \\
  \Se(\mathcal{I}_i) &= \text{normalised } \mu_i.
\end{align}
For a network at contact, the conductances and chemical potentials are
self-consistent: the flux pattern is the one that minimises total partition depth
subject to the stoichiometric constraints. This self-consistency produces maximal
correlation among the three coordinates (because all three are monotonically
related to the ``importance'' of each species in the network).

A perturbation $\delta$ modifies the effective conductance of one or more edges,
breaking the self-consistency. The perturbed flux pattern is no longer optimal
with respect to the contact conductances, introducing discrepancies between
$\Sk$ (which responds to flux changes) and $\St$ (which depends on
conductance, not flux) and $\Se$ (which depends on chemical potential, not
conductance). These discrepancies reduce the rank correlation, hence $R$
decreases. \qed
\end{proof}


%%%%%%%%%%%%%%%%%%%%%%%%%%%%%%%%%%%%%%%%%%%%%%%%%%%%%%%%%%%%%%%%%%%%%%%%%%%%%
\section{Contact Map Construction}
\label{sec:construction}
%%%%%%%%%%%%%%%%%%%%%%%%%%%%%%%%%%%%%%%%%%%%%%%%%%%%%%%%%%%%%%%%%%%%%%%%%%%%%

We now describe how to construct the contact map of a biochemical network from
publicly available pathway data.

\subsection{Input Data}

The construction requires three types of data for each species and reaction in
the network:

\begin{enumerate}[label=(\alph*)]
\item \textbf{Thermodynamic data}: standard Gibbs free energies of formation
  $\mu_i^\circ$ and reference concentrations $[\mathcal{I}_i]$ for each species.
  Sources: HMDB \cite{wishart2018hmdb}, KEGG \cite{kanehisa2000kegg}.

\item \textbf{Kinetic data}: catalytic rate constants $k_{\text{cat}}$ for each
  reaction. Sources: BRENDA \cite{schomburg2004brenda}, Reactome
  \cite{jassal2020reactome}.

\item \textbf{Network topology}: the set of reactions and their stoichiometry.
  Sources: KEGG \cite{kanehisa2000kegg}, Reactome \cite{jassal2020reactome},
  SBML models \cite{hucka2003sbml}.
\end{enumerate}

\subsection{Chemical Potential Computation}

For each species $\mathcal{I}_i$ with standard free energy of formation
$\mu_i^\circ$ (in kJ/mol) and concentration $[\mathcal{I}_i]$ (in M), the
chemical potential is:
\begin{equation}\label{eq:mu}
  \mu_i = \mu_i^\circ + RT \ln[\mathcal{I}_i],
\end{equation}
where $R = 8.314 \times 10^{-3}$ kJ/(mol$\cdot$K) is the gas constant and
$T = 310$ K is the physiological temperature \cite{atkins2014physical,
alberty2003thermodynamics}.

\subsection{Conductance Computation}

For each reaction $\mathcal{I}_i \to \mathcal{I}_j$ catalysed by an enzyme with
catalytic rate constant $k_{\text{cat}}$ (in s$^{-1}$), the conductance is
\cite{beard2002energy, noor2014pathway}:
\begin{equation}\label{eq:G}
  G_{ij} = \frac{k_{\text{cat}} \cdot [\mathcal{I}_i]}{RT},
\end{equation}
where $RT = 2.577$ kJ/mol at 310 K. This definition ensures that the flux
$J_{ij} = G_{ij}(\mu_i - \mu_j)$ has the correct units (mol/s per unit
volume) and is consistent with the Onsager reciprocal relations
\cite{onsager1931reciprocal}.

\subsection{S-Entropy Coordinate Computation}

Given $\{\mu_i\}$, $\{G_{ij}\}$, and the network topology $G = (V, E)$:

\begin{enumerate}[label=\arabic*.]
\item Compute the flux on each edge: $J_{ij} = G_{ij}(\mu_i - \mu_j)$.
\item Compute $\Sk(\mathcal{I}_i)$ via Equation~\eqref{eq:Sk}.
\item Compute $\St(\mathcal{I}_i)$ via Equation~\eqref{eq:St}.
\item Compute $\Se(\mathcal{I}_i)$ via Equation~\eqref{eq:Se}.
\item Compute the contact cost on each edge: $\CC(\mathcal{I}_i, \mathcal{I}_j)
  = \dS(\bS(\mathcal{I}_i), \bS(\mathcal{I}_j))$ via Equation~\eqref{eq:CM}.
\end{enumerate}

\begin{algorithm}[t]
\caption{Contact Map Construction}
\label{alg:construction}
\DontPrintSemicolon
\KwIn{Species set $\Omega$; reaction network $G = (V, E)$; thermodynamic data
  $\{\mu_i^\circ, [\mathcal{I}_i]\}$; kinetic data $\{k_{\text{cat},e}\}$}
\KwOut{Contact map $\CC: E \to \reals_{>0}$; S-entropy coordinates
  $\{\bS(\mathcal{I}_i)\}$}

\BlankLine
\tcp{Step 1: Chemical potentials}
\ForEach{$\mathcal{I}_i \in \Omega$}{
  $\mu_i \leftarrow \mu_i^\circ + RT \ln[\mathcal{I}_i]$\;
}

\BlankLine
\tcp{Step 2: Conductances}
\ForEach{edge $(\mathcal{I}_i, \mathcal{I}_j) \in E$}{
  $G_{ij} \leftarrow k_{\text{cat}} \cdot [\mathcal{I}_i] / RT$\;
}

\BlankLine
\tcp{Step 3: Fluxes}
\ForEach{edge $(\mathcal{I}_i, \mathcal{I}_j) \in E$}{
  $J_{ij} \leftarrow G_{ij} \cdot (\mu_i - \mu_j)$\;
}

\BlankLine
\tcp{Step 4: S-entropy coordinates}
$F_{\max} \leftarrow \max_{k \in V} \sum_{j \in \mathcal{N}(k)} |J_{kj}|$\;
$C_{\max} \leftarrow \max_{k \in V} \sum_{j \in \mathcal{N}(k)} G_{kj}$\;
$\mu_{\min} \leftarrow \min_{k \in V} \mu_k$; \quad
$\mu_{\max} \leftarrow \max_{k \in V} \mu_k$\;

\ForEach{$\mathcal{I}_i \in \Omega$}{
  $\Sk(\mathcal{I}_i) \leftarrow \sum_{j \in \mathcal{N}(i)} |J_{ij}| / F_{\max}$\;
  $\St(\mathcal{I}_i) \leftarrow \sum_{j \in \mathcal{N}(i)} G_{ij} / C_{\max}$\;
  $\Se(\mathcal{I}_i) \leftarrow (\mu_i - \mu_{\min}) / (\mu_{\max} - \mu_{\min})$\;
}

\BlankLine
\tcp{Step 5: Contact costs}
\ForEach{edge $(\mathcal{I}_i, \mathcal{I}_j) \in E$}{
  $\CC(\mathcal{I}_i, \mathcal{I}_j) \leftarrow \dS(\bS(\mathcal{I}_i), \bS(\mathcal{I}_j))$\;
}

\Return{$\CC$, $\{\bS(\mathcal{I}_i)\}$}\;
\end{algorithm}

\begin{proposition}[Complexity of Contact Map Construction]
\label{prop:complexity}
Algorithm~\ref{alg:construction} runs in $O(|V| + |E|)$ time and $O(|V| + |E|)$
space.
\end{proposition}

\begin{proof}
Steps 1 and 4 iterate over $|V|$ species; Steps 2, 3, and 5 iterate over $|E|$
edges. The maximum computations in Step 4 require one pass over $|V|$ species.
All operations per species or edge are $O(1)$. Total: $O(|V| + |E|)$. \qed
\end{proof}

\begin{figure*}[!htbp]
    \centering
    \includegraphics[width=\textwidth]{figures/panel_2_contact_graphs.png}
    \caption{\textbf{Contact graphs $G_C = (V, E, w)$ for four biochemical networks.}
    Vertices are coloured by $\Se$ (viridis scale) and edges are coloured by contact cost $\mathcal{C}(e)$ (RdYlGn diverging scale: green = low cost / tight contact, red = high cost / weak contact), with edge width inversely proportional to $\mathcal{C}(e)$.
    (\textbf{A})~Glycolysis: the path graph $P_{10}$ reveals the G6P$\to$F6P edge as the thinnest (highest cost, $\mathcal{C} = 1.354$), identifying phosphoglucose isomerase as the principal bottleneck separating upper and lower glycolysis. The 3PG$\to$2PG edge is the thickest (lowest cost, $\mathcal{C} = 0.002$), indicating near-equilibrium operation.
    (\textbf{B})~TCA cycle: the cycle graph $C_9$ with circular layout; the Citrate$\to$Isocitrate and Isocitrate$\to$$\alpha$KG edges carry the highest costs, consistent with the known regulatory role of isocitrate dehydrogenase.
    (\textbf{C})~Oxidative phosphorylation: the mixed graph spans three compartments; cross-compartment edges (e.g.\ NADH$\to$CoQ) appear as longer-range connections with moderate cost.
    (\textbf{D})~EGFR/MAPK signalling: the directed graph with feedback edge (ERK$\to$EGFR) shows tight coupling (low cost) along the core cascade and higher cost at the receptor--adaptor interface.}
    \label{fig:panel-contact-graphs}
\end{figure*}


%%%%%%%%%%%%%%%%%%%%%%%%%%%%%%%%%%%%%%%%%%%%%%%%%%%%%%%%%%%%%%%%%%%%%%%%%%%%%
\section{Perturbation Theory on Contact Maps}
\label{sec:perturbation}
%%%%%%%%%%%%%%%%%%%%%%%%%%%%%%%%%%%%%%%%%%%%%%%%%%%%%%%%%%%%%%%%%%%%%%%%%%%%%

We develop the theory of perturbations to the contact map, formalising the
concepts of disease, drug response, and restoration.

\subsection{Perturbation as Partition Depth Inflation}

\begin{definition}[Network Perturbation]
\label{def:perturbation}
A \emph{perturbation} of a biochemical contact map $\CC$ is a function
$\delta: E \to \reals_{\geq 0}$ that modifies the effective conductance of each
edge:
\begin{equation}\label{eq:perturb}
  G'_{ij} = \alpha_{ij} \cdot G_{ij},
  \qquad \alpha_{ij} \in (0, \infty),
\end{equation}
where $\alpha_{ij} \neq 1$ indicates a perturbation. An \emph{inhibition}
corresponds to $\alpha_{ij} < 1$ (reduced conductance); an \emph{activation}
corresponds to $\alpha_{ij} > 1$ (increased conductance). The resulting
change in partition depth along edge $(\mathcal{I}_i, \mathcal{I}_j)$ is:
\begin{equation}\label{eq:depth_change}
  \Delta D_{ij} = |D'_{ij} - D_{ij}|
  = |\dS(\bS'(\mathcal{I}_i), \bS'(\mathcal{I}_j))
    - \dS(\bS(\mathcal{I}_i), \bS(\mathcal{I}_j))|.
\end{equation}
\end{definition}

\begin{remark}
Both inhibitions and activations increase partition depth relative to the
contact state: an inhibition raises the thermodynamic barrier by reducing
conductance; an activation raises it by creating a non-equilibrium excess flux
that was not present at contact. In both cases, the perturbed state has higher
total partition depth than the contact state.
\end{remark}

\begin{definition}[Flux Visibility]
\label{def:visibility}
The \emph{flux visibility} $V$ of a perturbed contact map relative to the
healthy (contact) state is the weighted geometric mean of the per-edge flux
ratios:
\begin{equation}\label{eq:visibility}
  V = \prod_{(i,j) \in E}
    \left(\frac{\min(|J'_{ij}|, |J_{ij}|)}{\max(|J'_{ij}|, |J_{ij}|)}\right)^{w_{ij}},
\end{equation}
where $w_{ij} = G_{ij} / \sum_{(k,l) \in E} G_{kl}$ is the normalised
conductance weight and the product is taken over all edges. $V \in [0, 1]$,
with $V = 1$ at the contact state (no perturbation) and $V \to 0$ as the
perturbation becomes severe.
\end{definition}

\begin{proposition}[Visibility Decreases Under Perturbation]
\label{prop:visibility_decrease}
For any perturbation $\alpha \neq \mathbf{1}$, the flux visibility satisfies
$V < 1$.
\end{proposition}

\begin{proof}
If $\alpha_{ij} \neq 1$ for any edge, then $J'_{ij} = G'_{ij}(\mu'_i -
\mu'_j) \neq J_{ij}$ (because $G'_{ij} \neq G_{ij}$ and the chemical
potentials adjust via mass balance). Therefore $\min(|J'_{ij}|, |J_{ij}|) /
\max(|J'_{ij}|, |J_{ij}|) < 1$ for at least one edge, and since
$w_{ij} > 0$ for all edges, the geometric mean is strictly less than 1. \qed
\end{proof}

\subsection{Backward Navigation}

\begin{definition}[Backward Navigation]
\label{def:backward}
Given a perturbed contact map with a designated \emph{target species}
$\mathcal{I}_T$, the \emph{backward navigation} is the sequence of species
$\mathcal{I}_T, \mathcal{I}_{k_1}, \mathcal{I}_{k_2}, \ldots$ obtained by
greedily following the edge of maximum conductance from each species:
\begin{equation}\label{eq:backward}
  \mathcal{I}_{k_{n+1}}
  = \arg\max_{j \in \mathcal{N}(k_n)} G_{k_n j}.
\end{equation}
This traces the path of maximum thermodynamic coupling backward through the
network, identifying the ``causal chain'' of the most influential species
upstream of the target.
\end{definition}

\begin{proposition}[Backward Navigation Terminates]
\label{prop:backward_terminates}
In a finite, connected network, backward navigation terminates in at most $|V|$
steps, either by reaching a source node (species with no incoming edges in the
directed flux graph) or by entering a cycle.
\end{proposition}

\begin{proof}
Each step moves to a distinct neighbour (since $\arg\max$ over a finite set is
well-defined). If no neighbour has been previously visited, the path extends to
a new node. Since $|V|$ is finite, the path either reaches a node with no
further neighbours (a source) or revisits a node (entering a cycle) within $|V|$
steps. \qed
\end{proof}

\subsection{$\ell^1$-Optimal Restoration}

\begin{definition}[$\ell^1$-Optimal Restoration]
\label{def:l1_restoration}
Given a perturbed contact map with perturbation $\alpha$, the
\emph{$\ell^1$-optimal restoration} is the sparsest perturbation $\alpha^*$ that
restores the flux visibility above a threshold $V_0$:
\begin{equation}\label{eq:l1_restore}
  \alpha^* = \arg\min_{\alpha'} \|\alpha' - \mathbf{1}\|_1
  \quad\text{subject to}\quad
  V(\alpha') \geq V_0,
\end{equation}
where $\|\cdot\|_1$ is the $\ell^1$ norm (sum of absolute deviations from 1)
and $V(\alpha')$ is the flux visibility under the composite perturbation
$\alpha' \circ \alpha^{-1}$ \cite{foucart2013mathematical, candes2006robust}.
\end{definition}

\begin{remark}
The $\ell^1$ norm promotes sparsity: the optimal restoration modifies the fewest
possible edges \cite{foucart2013mathematical}. This corresponds to the drug
design problem of finding the fewest molecular targets whose modulation restores
normal network function \cite{hopkins2008network, csermely2013structure}.
\end{remark}

\begin{proposition}[Existence of Restoration]
\label{prop:restoration_exists}
For any perturbed network with $V < V_0$ and any threshold $V_0 \in (0, 1)$,
there exists at least one restoration $\alpha^*$ achieving $V(\alpha^*) \geq
V_0$, provided the network is connected.
\end{proposition}

\begin{proof}
The trivial restoration $\alpha^* = \alpha^{-1}$ (undoing the perturbation
entirely) achieves $V = 1 \geq V_0$. Since $\|\alpha^{-1} - \mathbf{1}\|_1$
is finite (all $\alpha_{ij}$ are positive and finite), the feasible set of the
optimisation problem~\eqref{eq:l1_restore} is non-empty. The $\ell^1$ norm is
continuous and the feasible set is closed and bounded (since all $\alpha'_{ij}
\in (0, M]$ for some finite $M$), so by the extreme value theorem, the minimum
exists. \qed
\end{proof}

\begin{figure*}[!htbp]
    \centering
    \includegraphics[width=\textwidth]{figures/panel_3_flux_visibility.png}
    \caption{\textbf{Flux visibility $V$ under single-edge perturbation across four disease models.}
    (\textbf{A})~$V(\alpha)$ response curves on a logarithmic $\alpha$-axis for all four pathways. Glycolysis (blue) shows the steepest decline under inhibition ($\alpha < 1$), dropping to $V = 0.143$ at $\alpha = 0.1$ (HK1 deficiency), while TCA (orange), OxPhos (green), and EGFR/MAPK (red) maintain $V > 0.9$ across the same perturbation range, reflecting the greater topological robustness of cyclic and branched graphs relative to path graphs. All curves satisfy $V(\alpha = 1) = 1$ (contact state).
    (\textbf{B})~Three-dimensional $V$ landscape: each pathway's $V(\alpha)$ curve is rendered at a distinct $y$-position, revealing that glycolysis occupies a qualitatively different response surface from the other three pathways.
    (\textbf{C})~Glycolysis HK1 perturbation: bar chart of $V$ across the $\alpha$-sweep $\{0.01, 0.05, \ldots, 10.0\}$. The contact state ($\alpha = 1.0$, green bar) achieves $V = 1$; both inhibition ($\alpha < 1$, red/orange) and activation ($\alpha > 1$) reduce $V$, confirming the symmetric visibility loss predicted by the theory. Bars transition from red ($V < 0.5$) through orange to green at $\alpha = 1$.
    (\textbf{D})~Disease visibility comparison: horizontal bar chart of $V_{\text{diseased}}$ for all four disease models. HK1 deficiency ($V = 0.143$) produces the most severe disruption; KRAS G12V ($V \approx 1.0$) produces negligible global visibility change despite 5-fold gain-of-function, consistent with the signalling cascade absorbing gain perturbations locally.}
    \label{fig:panel-flux-visibility}
\end{figure*}


%%%%%%%%%%%%%%%%%%%%%%%%%%%%%%%%%%%%%%%%%%%%%%%%%%%%%%%%%%%%%%%%%%%%%%%%%%%%%
\section{Computational Validation}
\label{sec:applications}
%%%%%%%%%%%%%%%%%%%%%%%%%%%%%%%%%%%%%%%%%%%%%%%%%%%%%%%%%%%%%%%%%%%%%%%%%%%%%

We now express the contact map framework entirely in graph-theoretic terms and
validate each claim computationally. All quantities are computed from publicly
available data: standard Gibbs energies of formation from eQuilibrator
\cite{noor2014pathway} and Alberty \cite{alberty2003thermodynamics}; catalytic
rate constants $k_{\text{cat}}$ from BRENDA \cite{schomburg2004brenda} (human
enzymes, pH~7.4, 37$^\circ$C); physiological concentrations from HMDB
\cite{wishart2018hmdb}; and network topology from KEGG \cite{kanehisa2000kegg}
and Reactome \cite{jassal2020reactome}. The computation uses \emph{zero free
parameters}: every input is a measured or catalogued quantity.

\subsection{Graph-Theoretic Nomenclature}

In graph-theoretic terms, the objects of the theory are:
\begin{itemize}
\item A \emph{contact graph} $G_C = (V, E, w)$: a finite weighted undirected
  graph where $V$ is the vertex set (molecular species), $E$ is the edge set
  (enzymatic reactions), and $w: E \to \reals_{>0}$ is the edge weight function
  (irreducible residue $\beta$).
\item A \emph{vertex embedding} $\bS: V \to [0,1]^3$: each vertex is mapped to
  a point in the S-entropy state space.
\item The \emph{contact cost} $\CC(e) = \dS(\bS(u), \bS(v))$ for each edge
  $e = (u,v) \in E$: the Euclidean distance between endpoint embeddings.
\item The \emph{contact Laplacian} $L_C$: the graph Laplacian of $(V, E,
  w^{-1})$ where edge weights are inverse contact costs.
\item A \emph{perturbation} $\alpha: E \to (0, \infty)$: a multiplicative
  modification of edge conductances, with $\alpha(e) = 1$ denoting no change.
\item The \emph{contact filtration}: the nested sequence of subgraphs
  $G_{\tau_1} \subseteq G_{\tau_2} \subseteq \cdots$ obtained by including
  edges in order of increasing contact cost.
\end{itemize}

\subsection{Glycolysis: A Path Graph $P_{10}$}
\label{sec:glycolysis}

The glycolysis pathway is a path graph $P_{10}$ on ten vertices (Glucose,
G6P, F6P, FBP, G3P, BPG$_{1,3}$, 3PG, 2PG, PEP, Pyruvate) with nine edges.

\paragraph{Validated S-entropy coordinates.}
Algorithm~\ref{alg:construction} yields the vertex embedding:
\begin{center}
\begin{tabular}{lcccc}
\toprule
Vertex $v$ & $\Sk(v)$ & $\St(v)$ & $\Se(v)$ & $\mu$ (kJ/mol) \\
\midrule
Glucose    & 0.999 & 0.921 & 0.771 & $-930.7$ \\
G6P        & 1.000 & 1.000 & 0.554 & $-1342.2$ \\
F6P        & 0.006 & 0.081 & 0.550 & $-1349.5$ \\
FBP        & 0.005 & 0.002 & 0.086 & $-2228.8$ \\
G3P        & 0.006 & 0.002 & 0.569 & $-1313.0$ \\
BPG$_{1,3}$ & 0.006 & 0.002 & 0.000 & $-2391.6$ \\
3PG        & 0.001 & 0.074 & 0.457 & $-1525.3$ \\
2PG        & 0.001 & 0.075 & 0.458 & $-1523.8$ \\
PEP        & 0.022 & 0.012 & 0.578 & $-1296.5$ \\
Pyruvate   & 0.020 & 0.010 & 1.000 & $-497.5$ \\
\bottomrule
\end{tabular}
\end{center}

\paragraph{Contact filtration.}
The filtration ordering (edges sorted by $\CC$) reveals the hierarchical
structure of glycolysis:
\begin{center}
\begin{tabular}{lcc}
\toprule
Edge $e$ & $\CC(e)$ & Components after \\
\midrule
3PG $\to$ 2PG       & 0.0018 & 9 \\
2PG $\to$ PEP       & 0.1373 & 8 \\
Glucose $\to$ G6P   & 0.2312 & 7 \\
PEP $\to$ Pyruvate  & 0.4219 & 6 \\
BPG$_{1,3}$ $\to$ 3PG & 0.4629 & 5 \\
F6P $\to$ FBP       & 0.4708 & 4 \\
FBP $\to$ G3P       & 0.4835 & 3 \\
G3P $\to$ BPG$_{1,3}$ & 0.5694 & 2 \\
G6P $\to$ F6P       & 1.3540 & 1 \\
\bottomrule
\end{tabular}
\end{center}
The first merge (3PG--2PG at $\CC = 0.0018$) identifies the tightest contact
in glycolysis: phosphoglycerate mutase operates near equilibrium. The last
merge (G6P--F6P at $\CC = 1.354$) identifies the bottleneck: phosphoglucose
isomerase separates the upper and lower halves of the pathway.

\begin{figure*}[!htbp]
    \centering
    \includegraphics[width=\textwidth]{figures/panel_4_contact_filtration.png}
    \caption{\textbf{Contact filtration and persistent hierarchy of four biochemical networks.}
    (\textbf{A})~Glycolysis filtration barcode: horizontal bars show the contact cost $\mathcal{C}(e)$ at which each edge is added to the filtration. The first merge (3PG--2PG, $\mathcal{C} = 0.002$) identifies the tightest contact in glycolysis; the last merge (G6P--F6P, $\mathcal{C} = 1.354$) identifies the principal bottleneck separating upper and lower glycolysis. The ordering recovers the biochemical hierarchy without biological priors.
    (\textbf{B})~Filtration step functions: connected component count versus threshold $\tau$ for all four pathways. Glycolysis (blue) shows a sharp late drop at $\tau \approx 1.35$ (the G6P--F6P bottleneck); TCA (orange) merges more uniformly; EGFR/MAPK (red) shows a stepped descent reflecting its multi-compartment structure.
    (\textbf{C})~Three-dimensional persistence landscape: filtration step functions rendered at distinct $y$-positions per pathway. The $z$-axis (components) decreases monotonically with threshold $\tau$, confirming the monotone inclusion property $G_{\tau_1} \subseteq G_{\tau_2}$ for $\tau_1 < \tau_2$.
    (\textbf{D})~Edge weight distribution (violin plots): contact cost $\mathcal{C}(e)$ distributions for all four pathways. Glycolysis shows the widest spread (median $\approx 0.47$, range $0.002$--$1.354$); TCA Cycle has a bimodal distribution; OxPhos is concentrated near $\mathcal{C} \approx 0.5$; EGFR/MAPK shows the tightest distribution with the lowest median, consistent with strong coupling along the signalling cascade.}
    \label{fig:panel-contact-filtration}
\end{figure*}

\paragraph{Spectral bisection.}
The Fiedler vector $\mathbf{f}_2$ of the contact Laplacian $L_C$ (eigenvector
for $\lambda_2 = 0.197$) partitions glycolysis into:
\[
  V^+ = \{\text{Glucose, G6P, F6P, FBP}\}, \quad
  V^- = \{\text{G3P, BPG}_{1,3}, \text{3PG, 2PG, PEP, Pyruvate}\},
\]
corresponding to the preparatory phase (ATP investment) and the payoff phase
(ATP generation) --- the classical biochemical division, recovered purely from
the contact Laplacian spectrum without any biological prior.

\paragraph{Triple coherence.}
$R = 0.463$, with pairwise Spearman correlations $\rho(\Sk, \St) = 0.333$,
$\rho(\St, \Se) = 0.370$, $\rho(\Sk, \Se) = 0.685$. The moderate coherence
reflects the heterogeneous distribution of flux control along the pathway
\cite{kacser1973control, fell1992metabolic}.

\paragraph{Perturbation: hexokinase deficiency.}
Hexokinase (HK1, EC 2.7.1.1) deficiency is modelled as $\alpha = 0.1$ on the
Glucose $\to$ G6P edge. Validated results:
$V = 0.143$ (severe flux disruption). Backward navigation from Pyruvate
traverses the full path graph in reverse: Pyruvate $\to$ PEP $\to$ 2PG $\to$
3PG $\to$ BPG$_{1,3}$ $\to$ G3P $\to$ FBP $\to$ F6P $\to$ G6P $\to$
\textbf{Glucose} (the perturbed vertex), correctly identifying the upstream
bottleneck. The $\ell^1$-optimal restoration identifies a single target
(sparsity~$= 1$), consistent with hexokinase activation as a therapeutic
strategy \cite{csermely2013structure}.

\subsection{TCA Cycle: A Cycle Graph $C_9$}

The TCA cycle is a cycle graph $C_9$ on nine vertices (Acetyl-CoA, Citrate,
Isocitrate, $\alpha$-KG, Succinyl-CoA, Succinate, Fumarate, Malate, OAA) with
nine edges, all localised to the mitochondrial matrix (single compartment).
The contact Laplacian has algebraic connectivity $\lambda_2 = 0.627$ --- the
highest among the four pathways, reflecting the topological robustness of the
cyclic structure \cite{kitano2004biological}.

\paragraph{Perturbation: IDH2 R172K.}
Isocitrate dehydrogenase 2 (IDH2) R172K is a neomorphic mutation common in
low-grade gliomas \cite{csermely2013structure}. We model 85\% activity
reduction: $\alpha = 0.15$ on the Isocitrate $\to$ $\alpha$-KG edge.
Validated: $V = 0.984$. The cyclic topology distributes the perturbation,
yielding a higher visibility than the path graph under comparable inhibition.
Backward navigation from OAA traces the full cycle: OAA $\to$ Malate $\to$
Fumarate $\to$ Succinate $\to$ Succinyl-CoA $\to$ $\alpha$-KG $\to$
\textbf{Isocitrate} $\to$ Citrate $\to$ Acetyl-CoA, passing through the
perturbed vertex.

\subsection{Oxidative Phosphorylation: A Mixed Graph $G_{\text{OxPhos}}$}

Oxidative phosphorylation is a mixed graph $G_{\text{OxPhos}}$ on eight vertices
spanning three compartments (mitochondrial matrix, inner mitochondrial membrane,
intermembrane space) with seven edges including cross-compartment coupling.
Algebraic connectivity: $\lambda_2 = 0.244$.

\paragraph{Compartmental factorisation}
The spectral bisection separates the graph into:
\[
  V^+ = \{\text{NADH, ADP, P}_i, \text{ATP}\}, \quad
  V^- = \{\text{CoQ, CytC, O}_2, \text{H}_2\text{O}\},
\]
which corresponds exactly to the substrate/product division: $V^+$ contains the
matrix-resident energy carriers, $V^-$ contains the electron transport chain
intermediates.

\paragraph{Perturbation: Complex I inhibition (rotenone).}
70\% inhibition ($\alpha = 0.3$) on NADH $\to$ CoQ. Validated: $V = 0.936$.
Backward navigation from ATP: ATP $\to$ ADP $\to$ \textbf{NADH} $\to$ CoQ
$\to$ CytC $\to$ O$_2$ $\to$ H$_2$O, identifying NADH as the upstream
bottleneck within three steps.

\subsection{EGFR/MAPK Signalling: A Directed Graph $G_{\text{MAPK}}$}

EGFR/MAPK signalling is a directed graph $G_{\text{MAPK}}$ on ten vertices
spanning four compartments (extracellular, membrane, cytoplasm, nucleus) with
ten edges including one feedback edge (ERK $\to$ EGFR). Algebraic connectivity:
$\lambda_2 = 0.705$, the highest spectral gap, reflecting the strong coupling
along the signalling cascade.

\paragraph{Compartmental factorisation.}
Validated: mean intra-compartmental contact cost $= 0.070$; mean
inter-compartmental contact cost $= 0.622$. The ratio
$\bar{\CC}_{\text{inter}} / \bar{\CC}_{\text{intra}} = 8.9$ confirms
Proposition~\ref{prop:factorisation}: membrane-crossing edges carry
substantially higher partition depth.

\paragraph{Perturbation: KRAS G12V.}
Constitutive activation modelled as $\alpha = 5.0$ on SOS $\to$ RAS.
Validated: $V = 0.9999$. The near-unity visibility indicates that the
signalling cascade, being a low-concentration system, absorbs gain-of-function
perturbations with minimal flux disruption at the global level. The
perturbation is detected not through $V$ but through localised S-entropy
shifts: the SOS $\to$ RAS edge moves furthest from contact.
Backward navigation from CycD traces the full cascade: CycD $\to$ MYC $\to$
ERK $\to$ MEK $\to$ RAF $\to$ RAS $\to$ \textbf{SOS} $\to$ GRB2 $\to$
EGFR $\to$ EGF.

\subsection{Cross-Pathway Summary}

\begin{center}
\begin{tabular}{lccccc}
\toprule
Pathway & $|V|$ & $|E|$ & $\lambda_2$ & $R$ & Graph type \\
\midrule
Glycolysis  & 10 & 9  & 0.197 & 0.463 & Path $P_{10}$ \\
TCA Cycle   & 9  & 9  & 0.627 & 0.111 & Cycle $C_9$ \\
OxPhos      & 8  & 7  & 0.244 & $-0.079$ & Mixed \\
EGFR/MAPK   & 10 & 10 & 0.705 & 0.152 & Directed + feedback \\
\bottomrule
\end{tabular}
\end{center}

\begin{center}
\begin{tabular}{lcccc}
\toprule
Disease model & $\alpha$ & $V$ & Nav.\ length & Sparsity \\
\midrule
HK1 deficiency     & 0.10 & 0.143 & 10 & 1 \\
IDH2 R172K         & 0.15 & 0.984 &  9 & 1 \\
Complex I rotenone  & 0.30 & 0.936 &  7 & 1 \\
KRAS G12V           & 5.00 & 1.000 & 10 & 1 \\
\bottomrule
\end{tabular}
\end{center}

The cross-pathway comparison reveals a structural insight: the algebraic
connectivity $\lambda_2$ of the contact Laplacian correlates with the graph's
topological complexity. The cycle graph $C_9$ (TCA) has higher $\lambda_2$
than the path graph $P_{10}$ (glycolysis), consistent with the general result
that cycles are more robust to edge removal than paths
\cite{chung1997spectral}.

\begin{figure*}[!htbp]
    \centering
    \includegraphics[width=\textwidth]{figures/panel_6_disease_profiles.png}
    \caption{\textbf{Disease perturbation profiles: S-entropy shifts, backward navigation, and restoration efficacy.}
    (\textbf{A})~Three-dimensional S-entropy shift vectors for hexokinase deficiency (glycolysis, $\alpha = 0.1$). Green markers denote healthy vertex positions $\mathbf{S}^{(0)}$; red markers denote diseased positions $\mathbf{S}^{(\alpha)}$; grey connecting lines show the displacement. The largest shifts occur at F6P ($|\Delta\mathbf{S}| = 0.395$) and Glucose ($|\Delta\mathbf{S}| = 0.383$), the two vertices flanking the perturbed edge.
    (\textbf{B})~Per-species S-entropy shift magnitudes $|\Delta\mathbf{S}|$ for all four disease models (grouped bars). Glycolysis HK1 deficiency (blue) produces large shifts across multiple vertices; TCA IDH2 R172K (orange) localises shifts to Isocitrate and $\alpha$-KG ($|\Delta\mathbf{S}| = 0.041$ each); OxPhos rotenone (green) affects only NADH and CoQ; EGFR KRAS G12V (red) produces negligible shifts ($|\Delta\mathbf{S}| < 0.001$), consistent with $V \approx 1$.
    (\textbf{C})~Backward navigation paths: each curve traces the greedy max-conductance trajectory from the terminal vertex to the source. Star markers indicate the identified source vertex. Navigation length ranges from 7 steps (OxPhos) to 10 steps (glycolysis, EGFR/MAPK), with all four navigations correctly terminating at or passing through the perturbed vertex.
    (\textbf{D})~Restoration efficacy: grouped bar chart comparing $V_{\text{healthy}} = 1$ (green), $V_{\text{diseased}}$ (red), and $V_{\text{restored}}$ (blue) for each pathway. The $\ell^1$-optimal restoration achieves $V_{\text{restored}} = 1.0$ in all four cases with single-edge correction (sparsity $= 1$), confirming that each disease model admits a minimal restoration target coinciding with the perturbed edge.}
    \label{fig:panel-disease-profiles}
\end{figure*}


%%%%%%%%%%%%%%%%%%%%%%%%%%%%%%%%%%%%%%%%%%%%%%%%%%%%%%%%%%%%%%%%%%%%%%%%%%%%%
\section{Additional Theoretical Properties}
\label{sec:properties}
%%%%%%%%%%%%%%%%%%%%%%%%%%%%%%%%%%%%%%%%%%%%%%%%%%%%%%%%%%%%%%%%%%%%%%%%%%%%%

\subsection{Stability Under Concentration Perturbation}

\begin{proposition}[Lipschitz Stability of Contact Map]
\label{prop:stability}
Let $[\mathcal{I}]$ and $[\mathcal{I}] + \eps \xi$ be two concentration vectors,
with $\|\xi\|_2 = 1$ and $\eps > 0$ small. Then for any edge
$(\mathcal{I}_i, \mathcal{I}_j) \in E$:
\begin{equation}
  |\CC_\eps(\mathcal{I}_i, \mathcal{I}_j)
  - \CC(\mathcal{I}_i, \mathcal{I}_j)|
  \leq L \cdot \eps,
\end{equation}
where $L > 0$ is a Lipschitz constant depending on the network topology and
thermodynamic parameters.
\end{proposition}

\begin{proof}
The contact map $\CC = \dS(\bS_i, \bS_j)$ is a composition of smooth functions:
$\bS_i$ depends on $\mu_i = \mu_i^\circ + RT \ln[\mathcal{I}_i]$ (smooth in
$[\mathcal{I}_i] > 0$), on $G_{ij} = k_{\text{cat}} [\mathcal{I}_i] / RT$
(linear in $[\mathcal{I}_i]$), and on $J_{ij} = G_{ij}(\mu_i - \mu_j)$ (smooth
in both). The Euclidean distance $\dS$ is Lipschitz with constant 1. By the
chain rule, the composition is Lipschitz with constant $L = \max(\partial \bS /
\partial [\mathcal{I}])$, which is bounded for $[\mathcal{I}_i]$ bounded away
from 0 and $\infty$. \qed
\end{proof}

\subsection{Compartmental Structure}

\begin{definition}[Compartmental Contact Map]
\label{def:compartmental}
A \emph{compartmentalised contact map} is a contact map $\CC$ together with a
partition of the species set $\Omega = \bigsqcup_{c=1}^C \Omega_c$ into
compartments (e.g., cytoplasm, mitochondrial matrix, nucleus, membrane).
The contact map restricted to each compartment $\CC|_{\Omega_c}$ is the
\emph{intra-compartmental contact map}, and the edges between compartments
form the \emph{inter-compartmental contact map}.
\end{definition}

\begin{proposition}[Compartmental Factorisation]
\label{prop:factorisation}
In a compartmentalised contact map, the partition depth between two species in
different compartments is always greater than the partition depth between species
in the same compartment:
\begin{equation}
  D(\mathcal{I}_i, \mathcal{I}_j) > D(\mathcal{I}_k, \mathcal{I}_l)
\end{equation}
whenever $\mathcal{I}_i \in \Omega_c$, $\mathcal{I}_j \in \Omega_d$ with $c
\neq d$, and $\mathcal{I}_k, \mathcal{I}_l \in \Omega_c$ are both in the same
compartment.
\end{proposition}

\begin{proof}
Inter-compartmental transport requires crossing a membrane, which imposes an
additional thermodynamic barrier (the free energy of membrane translocation)
on top of the chemical potential difference between the species
\cite{diekmann2013evolution, zecchin2015metabolic}. This additional barrier
contributes positively to the partition depth, ensuring
$D_{\text{inter}} > D_{\text{intra}}$. \qed
\end{proof}

\subsection{Hierarchical Contact Resolution}

\begin{definition}[Contact Resolution Ordering]
\label{def:resolution_ordering}
The \emph{contact resolution ordering} is the sequence of edges
$(e_1, e_2, \ldots, e_{|E|})$ sorted by increasing contact cost:
$\CC(e_1) \leq \CC(e_2) \leq \cdots \leq \CC(e_{|E|})$. This ordering
defines a filtration of the contact graph: at threshold $\tau$, only edges
with $\CC(e) \leq \tau$ are ``resolved.''
\end{definition}

\begin{theorem}[Contact Filtration Produces a Persistent Hierarchy]
\label{thm:filtration}
The contact resolution ordering defines a nested sequence of subgraphs
$G_1 \subseteq G_2 \subseteq \cdots \subseteq G_{|E|}$ such that the
connected components of $G_t$ form a hierarchical partition of the species set.
This hierarchy is persistent: components that merge at threshold $\tau$ remain
merged at all higher thresholds $\tau' > \tau$ (by Contact Invariance,
Theorem~\ref{thm:contact_invariance}).
\end{theorem}

\begin{proof}
At each threshold $\tau$, the subgraph $G_\tau = (V, \{e \in E : \CC(e) \leq
\tau\})$ includes all edges with contact cost at most $\tau$. As $\tau$
increases, edges are added but never removed (monotonicity of inclusion).
When a new edge connects two previously disconnected components, those
components merge. By Contact Invariance, the merged component cannot split at
any higher threshold: the contact relationship is preserved under refinement.

This defines a dendrogram (hierarchical clustering) on the species set, where
the merge height of two species is their contact cost $\CC(\mathcal{I}_i,
\mathcal{I}_j)$. The dendrogram is well-defined and unique (up to the ordering
of species within each level). The persistence of the hierarchy follows from
the monotonicity of the contact record \cite{edelsbrunner2008persistent,
ghrist2008barcodes, carlsson2009topology}. \qed
\end{proof}

\begin{remark}
The contact filtration is closely related to persistent homology
\cite{edelsbrunner2008persistent, carlsson2009topology}: the connected
components at each threshold are the 0-dimensional persistent homology groups
of the Vietoris-Rips complex built on the S-entropy state space with distance
$\dS$. Higher-dimensional homology groups (cycles, cavities) capture more
refined topological features of the contact map, such as metabolic cycles
(1-dimensional) and compartmental enclosures (2-dimensional). We leave the
systematic application of persistent homology to contact maps for future work.
\end{remark}

\subsection{Spectral Properties}

\begin{definition}[Contact Laplacian]
\label{def:laplacian}
The \emph{contact Laplacian} is the matrix $L_C \in \reals^{N \times N}$
defined by:
\begin{equation}\label{eq:laplacian}
  (L_C)_{ij} =
  \begin{cases}
    -\CC(\mathcal{I}_i, \mathcal{I}_j)^{-1}
      & \text{if } (\mathcal{I}_i, \mathcal{I}_j) \in E, \\
    \sum_{k \neq i} \CC(\mathcal{I}_i, \mathcal{I}_k)^{-1}
      & \text{if } i = j, \\
    0 & \text{otherwise}.
  \end{cases}
\end{equation}
The inverse contact cost $\CC^{-1}$ serves as the edge weight: species in close
contact (low $\CC$) have high coupling in the Laplacian.
\end{definition}

\begin{proposition}[Spectral Gap and Robustness]
\label{prop:spectral_gap}
The algebraic connectivity $\lambda_2(L_C)$ (the second-smallest eigenvalue
of $L_C$) is a lower bound on the robustness of the contact map to edge
removal:
\begin{equation}
  \lambda_2(L_C)
  \leq \min_{S \subset V, 1 \leq |S| \leq |V|/2}
    \frac{\sum_{i \in S, j \notin S} \CC(\mathcal{I}_i, \mathcal{I}_j)^{-1}}
         {|S|}
\end{equation}
(the Cheeger inequality for graphs \cite{chung1997spectral}).
A large $\lambda_2$ indicates that the network cannot be easily disconnected
by removing a few edges, corresponding to biological robustness
\cite{albert2000error, kitano2004biological, stelling2004robustness}.
\end{proposition}

\begin{proof}
This is a direct application of the discrete Cheeger inequality
\cite{chung1997spectral} to the weighted graph $(V, E, \CC^{-1})$. \qed
\end{proof}

\begin{figure*}[!htbp]
    \centering
    \includegraphics[width=\textwidth]{figures/panel_5_spectral_properties.png}
    \caption{\textbf{Contact Laplacian spectral properties for four biochemical networks.}
    (\textbf{A})~Eigenvalue spectra $\{\lambda_i\}$ of the contact Laplacian $L_C$ (symlog scale). All four pathways satisfy $\lambda_1 \approx 0$ (connected graph) and $\lambda_2 > 0$ (algebraic connectivity). The eigenvalue spread reflects graph complexity: EGFR/MAPK (red) and TCA (orange) show steeper spectral ramps than glycolysis (blue), consistent with their richer connectivity.
    (\textbf{B})~Fiedler vector heatmap: the second eigenvector $\mathbf{f}_2$ of $L_C$ for each pathway, displayed as a colour matrix (RdBu diverging scale). Sign changes in $\mathbf{f}_2$ identify the spectral bisection point. Glycolysis shows a single clean sign change between vertex indices 3 and 4, separating the preparatory phase (Glucose--FBP) from the payoff phase (G3P--Pyruvate).
    (\textbf{C})~Three-dimensional spectral bisection of glycolysis: vertices plotted in $(\Sk, \St, f_2)$ space, coloured by Fiedler sign (blue = positive, red = negative). The spatial separation confirms that the spectral bisection recovers the classical biochemical division of glycolysis into ATP-investment and ATP-generation phases purely from the contact Laplacian, without biological annotation.
    (\textbf{D})~Algebraic connectivity $\lambda_2$ comparison: EGFR/MAPK achieves the highest $\lambda_2 = 0.705$, followed by TCA ($\lambda_2 = 0.627$), OxPhos ($\lambda_2 = 0.244$), and glycolysis ($\lambda_2 = 0.197$). The ranking reflects topological robustness: cyclic and feedback-containing graphs resist edge removal better than path graphs, consistent with the general spectral graph theory result $\lambda_2(C_n) > \lambda_2(P_n)$.}
    \label{fig:panel-spectral-properties}
\end{figure*}


%%%%%%%%%%%%%%%%%%%%%%%%%%%%%%%%%%%%%%%%%%%%%%%%%%%%%%%%%%%%%%%%%%%%%%%%%%%%%
\section{Discussion}
\label{sec:discussion}
%%%%%%%%%%%%%%%%%%%%%%%%%%%%%%%%%%%%%%%%%%%%%%%%%%%%%%%%%%%%%%%%%%%%%%%%%%%%%

\subsection{The Contact State as Privileged Reference}

The central claim of this paper is that the contact map provides a unique
privileged reference state for any biochemical network. Unlike steady-state
models (which assume a particular flux distribution) or kinetic models
(which require time-dependent simulation), the contact map captures the
\emph{invariant skeleton} of the network: the set of minimum thermodynamic
separations that cannot be altered by any physical perturbation.

This has immediate practical consequences:

\begin{enumerate}[label=(\roman*)]
\item \textbf{Disease characterisation.} A disease state is a departure from
  contact along specific edges. The triple coherence $R$ provides a
  single-number summary of how far the network has departed from its contact
  state. This is a computable, parameter-free diagnostic quantity that can be
  derived from standard metabolomics and fluxomics data.

\item \textbf{Drug target identification.} The $\ell^1$-optimal restoration
  identifies the sparsest set of edges whose modulation restores the network
  toward contact. This directly addresses the polypharmacology question
  \cite{hopkins2008network}: which combination of targets is most efficient?

\item \textbf{Cross-network comparison.} Two networks with isomorphic contact
  maps (Theorem~\ref{thm:contact_completeness}) have the same partition
  topology, regardless of differences in absolute concentrations, flux
  magnitudes, or compartment sizes. This enables principled comparison of
  networks across cell types, organisms, or experimental conditions.
\end{enumerate}

\subsection{Relationship to Existing Frameworks}

The contact map framework differs from and complements existing approaches:

\paragraph{Flux Balance Analysis \cite{orth2010flux}.}
FBA finds a feasible flux vector satisfying stoichiometric constraints and
optimising an objective function. The contact map does not require an objective
function; instead, the contact state serves as the natural reference (the
partition ground state). FBA identifies \emph{what} fluxes are feasible; the
contact map identifies \emph{which flux distributions are close to or far from
the invariant structure of the network}.

\paragraph{Metabolic Control Analysis \cite{kacser1973control, fell1992metabolic,
heinrich1974linear}.}
MCA quantifies the sensitivity of steady-state fluxes and concentrations to
small perturbations of enzyme activities. The contact map provides a
complementary perspective: instead of linearising around a steady state, it
identifies the topological invariant (contact) and measures departures from it.
MCA control coefficients are local derivatives; contact costs are global
topological quantities.

\paragraph{Network topology \cite{barabasi2004network, jeong2000large,
ravasz2002hierarchical}.}
Pure topology analysis uses unweighted or uniformly weighted graphs. The contact
map is a \emph{thermodynamically weighted} topology: the edge weights are not
arbitrary but are determined by the irreducible residues $\beta_{ij}$, which
are intrinsic to the chemistry of the species pair. This grounding in
thermodynamics is what makes the contact map invariant under perturbation,
unlike arbitrary graph metrics.

\paragraph{Persistent homology \cite{carlsson2009topology,
edelsbrunner2008persistent, rabadan2020topological}.}
The contact filtration (Theorem~\ref{thm:filtration}) produces a persistent
hierarchy that is formally equivalent to a persistence diagram. The contact
map can therefore be analysed using the full machinery of topological data
analysis, including stability theorems and statistical tests. The novel
contribution here is the identification of the filtration parameter (S-entropy
distance) with a thermodynamic quantity (partition depth), providing a physical
interpretation of the persistence diagram.

\subsection{Limitations}

Several limitations deserve acknowledgment:

\begin{enumerate}[label=(\roman*)]
\item \textbf{Steady-state assumption.} The S-entropy coordinates are defined
  for a system at (quasi-)steady state. Extension to time-varying systems would
  require tracking the evolution of $\bS(\mathcal{I}_i, t)$ and defining
  contact as the infimum over time, which may not be achieved at any finite time.

\item \textbf{Compartment boundaries.} The compartmental factorisation
  (Proposition~\ref{prop:factorisation}) assumes that membrane translocation
  costs are uniform within each compartment pair. In reality, transporter
  specificities create heterogeneous inter-compartmental barriers.

\item \textbf{Stochastic effects.} In low-copy-number regimes (e.g., signalling
  networks with fewer than $\sim$100 molecules per cell), the chemical potential
  is not well-defined in the classical sense. A stochastic extension using
  chemical master equations would be needed.

\item \textbf{Parameter uncertainty.} The S-entropy coordinates depend on
  $k_{\text{cat}}$ values and concentrations that may have significant
  measurement uncertainty. The Lipschitz stability
  (Proposition~\ref{prop:stability}) bounds the propagation of this uncertainty
  but does not eliminate it.
\end{enumerate}


%%%%%%%%%%%%%%%%%%%%%%%%%%%%%%%%%%%%%%%%%%%%%%%%%%%%%%%%%%%%%%%%%%%%%%%%%%%%%
\section{Conclusion}
\label{sec:conclusion}
%%%%%%%%%%%%%%%%%%%%%%%%%%%%%%%%%%%%%%%%%%%%%%%%%%%%%%%%%%%%%%%%%%%%%%%%%%%%%

We have introduced a rigorous mathematical framework for representing
biochemical networks as contact maps: weighted graphs in which edge weights are
the irreducible thermodynamic residues of species individuation.

\subsection{Summary of Results}

The main results of this paper are:

\begin{enumerate}[label=(\roman*)]
\item \textbf{Axiomatic foundation} (Section~\ref{sec:foundations}): three
  postulates --- individuation by negation, finite individuation cost, and
  monotonicity of partition count --- suffice to generate the entire theory.
  Theorem~\ref{thm:noninstantaneous} establishes that species individuation
  requires positive thermodynamic cost.

\item \textbf{Contact invariance} (Theorem~\ref{thm:contact_invariance}):
  contact relationships are preserved under all physically realisable
  perturbations. The contact map is the invariant skeleton of the network.

\item \textbf{Contact irreducibility} (Theorem~\ref{thm:contact_irreducibility}):
  no coarser topological fact carries the same information as contact. Contact
  is the minimum sufficient invariant.

\item \textbf{Contact completeness} (Theorem~\ref{thm:contact_completeness}):
  the contact map is a complete invariant of the network's partition topology.
  Two networks with isomorphic contact maps have identical partition structure.

\item \textbf{S-entropy state space} (Section~\ref{sec:sentropy}): three
  independent entropy coordinates $(\Sk, \St, \Se)$ provide a computable,
  well-defined embedding of each species into $[0,1]^3$. The S-entropy distance
  is a metric (Theorem~\ref{thm:metric}) and is Lipschitz-stable under
  concentration perturbation (Proposition~\ref{prop:stability}).

\item \textbf{Contact filtration} (Theorem~\ref{thm:filtration}): the contact
  resolution ordering produces a persistent hierarchy on the species set,
  connecting contact maps to the theory of persistent homology.

\item \textbf{Applications}: explicit computation of contact maps, perturbation
  responses, and restoration targets for glycolysis, TCA cycle, oxidative
  phosphorylation, and EGFR/MAPK signalling, using zero free parameters.
\end{enumerate}

\subsection{The Contact Map as the State of a Cell}

The contact map answers the question posed in the introduction: \emph{what is
invariant?} When a cell is perturbed, its concentrations change, its fluxes
reroute, its S-entropy coordinates shift. But its contact map remains. The
disease is not a different cell; it is the same contact map with elevated
partition depth along specific edges. The healthy state is not one particular
concentration profile; it is the state in which all edges are at their minimum
partition depth --- the contact state.

This perspective inverts the traditional modelling paradigm: instead of
starting from a kinetic model and asking ``what happens when we perturb it?'',
we start from the contact map and ask ``how far has this state departed from
contact?'' The answer is a computable quantity (the triple coherence $R$, the
flux visibility $V$, or the $\ell^1$-restoration cost), and the prescription
for restoration is the $\ell^1$-optimal perturbation that moves the state back
toward contact.

The contact map is, in this precise sense, the state of a cell from which all
other states can be derived.

\printbibliography

\end{document}
